\pdfoutput=1
% ---------------------------------------------------------------------------
% Does coupling divergence tell a god unit from a complex algorithm?
% A pre-registered test of weighting complexity concentration by the data a
% unit hands on.
%
% Built for arXiv (primary class cs.SE). No .bib: the bibliography is a
% `thebibliography` environment at the end.
%
% Every number in this file is a macro from generated/numbers.tex and every
% table is a fragment under generated/, both written by paper/scripts/tables.py
% from the harness's confirmatory analysis. No number this study computes is
% typed by hand; the only typed numbers are literature constants and the
% plan's own constants. Code, commits, files and commands are named in the
% ancillary README, not here.
%
% Build:   pdflatex main && pdflatex main
% ---------------------------------------------------------------------------
\documentclass[11pt,a4paper]{article}

\usepackage[T1]{fontenc}
\usepackage[utf8]{inputenc}
\usepackage{lmodern}
\usepackage[margin=2.4cm]{geometry}
\usepackage{graphicx}
\usepackage{booktabs}
\usepackage{multirow}
\usepackage{amsmath}
\usepackage{microtype}
\usepackage[hidelinks]{hyperref}
\usepackage{caption}
\captionsetup{font=small,labelfont=bf}

% generated by paper/scripts/tables.py; do not edit
\newcommand{\analysisMadeUtc}{2026-10-26 14:55:06 UTC}
\newcommand{\analysisVOneMadeUtc}{2026-10-26 08:04:12 UTC}
\newcommand{\analysisKind}{confirmatory}
\newcommand{\analysisShaShort}{9aa901ddc12b}
\newcommand{\analysisVOneShaShort}{d64d4c8d547f}
\newcommand{\decisionsMadeUtc}{2026-10-26 13:53:46 UTC}
\newcommand{\bootB}{2{,}000}
\newcommand{\bootSeed}{12345}
\newcommand{\mlcqRepos}{522}
\newcommand{\mlcqSamples}{4{,}764}
\newcommand{\mlcqBlobReviewed}{2{,}334}
\newcommand{\mlcqBlob}{254}
\newcommand{\mlcqLongReviewed}{2{,}430}
\newcommand{\mlcqLong}{276}
\newcommand{\mlcqMovedMeasured}{23}
\newcommand{\mlcqMovedConfirm}{12}
\newcommand{\mlcqMovedWithBlob}{12}
\newcommand{\mlcqMeasured}{520}
\newcommand{\mlcqConfirmMeasured}{261}
\newcommand{\mlcqConfirmUsable}{258}
\newcommand{\mlcqConfirmErrors}{2}
\newcommand{\mlcqExploreMeasured}{261}
\newcommand{\seedThreeMeasured}{500}
\newcommand{\seedThreeUsable}{492}
\newcommand{\seedThreeDrawUsable}{580}
\newcommand{\seedThreeExcluded}{62}
\newcommand{\fontanaSystems}{73}
\newcommand{\fontanaInstances}{420}
\newcommand{\fontanaGodInstances}{140}
\newcommand{\fontanaFilesReviewed}{384}
\newcommand{\fontanaFilesGod}{139}
\newcommand{\fontanaFilesReviewedMeasured}{380}
\newcommand{\fontanaFilesGodMeasured}{139}
\newcommand{\validationRepos}{150}
\newcommand{\joinMlcqLizardUnits}{1{,}949{,}548}
\newcommand{\joinMlcqShare}{99.4\%}
\newcommand{\joinMlcqMeasurerExtra}{40{,}390}
\newcommand{\joinMlcqMeasurerExtraShare}{2.0\%}
\newcommand{\joinMlcqClassUnits}{179{,}585}
\newcommand{\joinMlcqClassCoupledShare}{99.9\%}
\newcommand{\callsMlcqInternal}{42.3\%}
\newcommand{\callsMlcqJdk}{25.8\%}
\newcommand{\callsMlcqUnresolved}{24.0\%}
\newcommand{\callsMlcqShouldHave}{7.9\%}
\newcommand{\callsMlcqBudget}{0.0\%}
\newcommand{\roleMlcqClassCatchmentUnits}{50.2\%}
\newcommand{\roleMlcqClassCatchmentCc}{45.5\%}
\newcommand{\roleMlcqClassPipeUnits}{6.6\%}
\newcommand{\roleMlcqClassPipeCc}{1.4\%}
\newcommand{\roleMlcqClassJunctionUnits}{43.2\%}
\newcommand{\roleMlcqClassJunctionCc}{53.1\%}
\newcommand{\roleMlcqMethodCatchmentUnits}{56.6\%}
\newcommand{\roleMlcqMethodCatchmentCc}{40.9\%}
\newcommand{\roleMlcqMethodPipeUnits}{12.4\%}
\newcommand{\roleMlcqMethodPipeCc}{9.1\%}
\newcommand{\roleMlcqMethodJunctionUnits}{31.0\%}
\newcommand{\roleMlcqMethodJunctionCc}{50.0\%}
\newcommand{\roleMlcqMethodDataCatchmentUnits}{35.2\%}
\newcommand{\roleMlcqMethodDataPipeUnits}{40.0\%}
\newcommand{\throwNewMlcqClassShare}{2.9\%}
\newcommand{\throwNewMlcqMethodShare}{2.1\%}
\newcommand{\diMlcqRepos}{113}
\newcommand{\prunedMlcqFiles}{1{,}271}
\newcommand{\filesMlcq}{486{,}712}
\newcommand{\filesMlcqFailed}{1{,}251}
\newcommand{\filesMlcqFailedShare}{0.26\%}
\newcommand{\scaleMlcqClassRepos}{259}
\newcommand{\scaleMlcqXMad}{257}
\newcommand{\scaleMlcqDMad}{94}
\newcommand{\scaleMlcqDMeanAd}{163}
\newcommand{\scaleMlcqDConstant}{2}
\newcommand{\zeroMlcqDClass}{57.3\%}
\newcommand{\scaleMlcqMethodRepos}{259}
\newcommand{\scaleMlcqXmMeanAd}{244}
\newcommand{\scaleMlcqDmMeanAd}{259}
\newcommand{\zeroMlcqXm}{63.5\%}
\newcommand{\zeroMlcqDmMethod}{80.5\%}
\newcommand{\joinSeedThreeLizardUnits}{2{,}993{,}414}
\newcommand{\joinSeedThreeShare}{98.3\%}
\newcommand{\joinSeedThreeMeasurerExtra}{89{,}757}
\newcommand{\joinSeedThreeMeasurerExtraShare}{3.0\%}
\newcommand{\joinSeedThreeClassUnits}{295{,}031}
\newcommand{\joinSeedThreeClassCoupledShare}{98.5\%}
\newcommand{\callsSeedThreeInternal}{37.7\%}
\newcommand{\callsSeedThreeJdk}{25.1\%}
\newcommand{\callsSeedThreeUnresolved}{27.3\%}
\newcommand{\callsSeedThreeShouldHave}{9.7\%}
\newcommand{\callsSeedThreeBudget}{0.0\%}
\newcommand{\roleSeedThreeClassCatchmentUnits}{50.7\%}
\newcommand{\roleSeedThreeClassCatchmentCc}{42.8\%}
\newcommand{\roleSeedThreeClassPipeUnits}{6.4\%}
\newcommand{\roleSeedThreeClassPipeCc}{1.5\%}
\newcommand{\roleSeedThreeClassJunctionUnits}{42.9\%}
\newcommand{\roleSeedThreeClassJunctionCc}{55.7\%}
\newcommand{\roleSeedThreeMethodCatchmentUnits}{57.0\%}
\newcommand{\roleSeedThreeMethodCatchmentCc}{41.5\%}
\newcommand{\roleSeedThreeMethodPipeUnits}{11.9\%}
\newcommand{\roleSeedThreeMethodPipeCc}{9.1\%}
\newcommand{\roleSeedThreeMethodJunctionUnits}{31.1\%}
\newcommand{\roleSeedThreeMethodJunctionCc}{49.4\%}
\newcommand{\roleSeedThreeMethodDataCatchmentUnits}{36.4\%}
\newcommand{\roleSeedThreeMethodDataPipeUnits}{38.6\%}
\newcommand{\throwNewSeedThreeClassShare}{2.2\%}
\newcommand{\throwNewSeedThreeMethodShare}{1.6\%}
\newcommand{\diSeedThreeRepos}{163}
\newcommand{\prunedSeedThreeFiles}{26{,}123}
\newcommand{\filesSeedThree}{359{,}093}
\newcommand{\filesSeedThreeFailed}{5{,}681}
\newcommand{\filesSeedThreeFailedShare}{1.58\%}
\newcommand{\scaleSeedThreeClassRepos}{500}
\newcommand{\scaleSeedThreeXMad}{495}
\newcommand{\scaleSeedThreeDMad}{122}
\newcommand{\scaleSeedThreeDMeanAd}{364}
\newcommand{\scaleSeedThreeDConstant}{14}
\newcommand{\zeroSeedThreeDClass}{65.6\%}
\newcommand{\scaleSeedThreeMethodRepos}{500}
\newcommand{\scaleSeedThreeXmMeanAd}{448}
\newcommand{\scaleSeedThreeDmMeanAd}{492}
\newcommand{\zeroSeedThreeXm}{64.9\%}
\newcommand{\zeroSeedThreeDmMethod}{83.5\%}
\newcommand{\joinFontanaLizardUnits}{376{,}803}
\newcommand{\joinFontanaShare}{99.3\%}
\newcommand{\joinFontanaMeasurerExtra}{4{,}594}
\newcommand{\joinFontanaMeasurerExtraShare}{1.2\%}
\newcommand{\joinFontanaClassUnits}{33{,}118}
\newcommand{\joinFontanaClassCoupledShare}{99.9\%}
\newcommand{\callsFontanaInternal}{49.9\%}
\newcommand{\callsFontanaJdk}{35.7\%}
\newcommand{\callsFontanaUnresolved}{11.1\%}
\newcommand{\callsFontanaShouldHave}{3.2\%}
\newcommand{\callsFontanaBudget}{0.0\%}
\newcommand{\roleFontanaClassCatchmentUnits}{37.9\%}
\newcommand{\roleFontanaClassCatchmentCc}{32.7\%}
\newcommand{\roleFontanaClassPipeUnits}{4.7\%}
\newcommand{\roleFontanaClassPipeCc}{1.1\%}
\newcommand{\roleFontanaClassJunctionUnits}{57.4\%}
\newcommand{\roleFontanaClassJunctionCc}{66.2\%}
\newcommand{\throwNewFontanaClassShare}{1.8\%}
\newcommand{\diFontanaRepos}{3}
\newcommand{\prunedFontanaFiles}{14}
\newcommand{\filesFontana}{37{,}781}
\newcommand{\filesFontanaFailed}{46}
\newcommand{\filesFontanaFailedShare}{0.12\%}
\newcommand{\scaleFontanaClassRepos}{73}
\newcommand{\scaleFontanaXMad}{73}
\newcommand{\scaleFontanaDMad}{54}
\newcommand{\scaleFontanaDMeanAd}{19}
\newcommand{\scaleFontanaDConstant}{0}
\newcommand{\zeroFontanaDClass}{40.0\%}
\newcommand{\uncoupledBlobReviewed}{934}
\newcommand{\uncoupledBlobUncoupled}{0}
\newcommand{\uncoupledLongReviewed}{1034}
\newcommand{\uncoupledLongUncoupled}{7}
\newcommand{\uncoupledLongPositiveUncoupled}{0}
\newcommand{\uncoupledFontanaUncoupled}{3}
\newcommand{\uncoupledFontanaPositiveUncoupled}{2}
\newcommand{\labelsBlobReviewedLabelled}{1{,}060}
\newcommand{\labelsBlobReviewedMeasured}{934}
\newcommand{\labelsBlobLabelled}{110}
\newcommand{\labelsBlobMeasured}{98}
\newcommand{\labelsBlobReviewedDeleted}{5}
\newcommand{\labelsLongReviewedLabelled}{1{,}170}
\newcommand{\labelsLongReviewedMeasured}{1{,}034}
\newcommand{\labelsLongLabelled}{134}
\newcommand{\labelsLongMeasured}{125}
\newcommand{\labelsLongReviewedDeleted}{2}
\newcommand{\mlcqOwnerTopShare}{50.8\%}
\newcommand{\seedThreeOwnerTopShare}{2.8\%}
\newcommand{\mlcqApacheRepos}{131}
\newcommand{\agreeCcExact}{99.7\%}
\newcommand{\agreeCcRho}{0.997}
\newcommand{\agreeCcRhoWithin}{0.998}
\newcommand{\agreeCognitiveExact}{94.6\%}
\newcommand{\agreeCognitiveRho}{0.959}
\newcommand{\agreeCognitiveRhoWithin}{0.972}
\newcommand{\agreeNlocExact}{98.3\%}
\newcommand{\agreeNlocRho}{0.998}
\newcommand{\agreeNlocRhoWithin}{0.999}
\newcommand{\agreeTokensExact}{94.6\%}
\newcommand{\agreeTokensRho}{0.998}
\newcommand{\agreeTokensRhoWithin}{0.999}
\newcommand{\agreeParamsExact}{99.9\%}
\newcommand{\agreeParamsRho}{0.999}
\newcommand{\agreeParamsRhoWithin}{1.000}
\newcommand{\agreeNestingExact}{93.3\%}
\newcommand{\agreeNestingRho}{0.911}
\newcommand{\agreeNestingRhoWithin}{0.939}
\newcommand{\agreeCycDensityExact}{98.3\%}
\newcommand{\agreeCycDensityRho}{0.997}
\newcommand{\agreeCycDensityRhoWithin}{0.998}
\newcommand{\agreeExitPointsExact}{99.1\%}
\newcommand{\agreeExitPointsRho}{0.989}
\newcommand{\agreeExitPointsRhoWithin}{0.994}
\newcommand{\agreeHvolExact}{1.5\%}
\newcommand{\agreeHvolRho}{0.990}
\newcommand{\agreeHvolRhoWithin}{0.990}
\newcommand{\agreeMiExact}{1.5\%}
\newcommand{\agreeMiRho}{0.995}
\newcommand{\agreeMiRhoWithin}{0.994}
\newcommand{\agreeUnits}{2{,}300{,}061}
\newcommand{\agreeRepos}{261}
\newcommand{\agreeJoinShare}{99.4\%}
\newcommand{\pmdCogPaperOne}{0.982}
\newcommand{\pmdCogMeasurer}{0.972}
\newcommand{\callsExploreInternalMedian}{34.8\%}
\newcommand{\callsExploreUnresolvedMedian}{26.9\%}
\newcommand{\callsExploreShouldHaveMedian}{6.8\%}
\newcommand{\xClass}{nloc\_max}
\newcommand{\xClassMcc}{0.287}
\newcommand{\xClassSecond}{hvol\_max}
\newcommand{\xClassSecondMcc}{0.279}
\newcommand{\xClassThird}{tokens\_max}
\newcommand{\xClassThirdMcc}{0.268}
\newcommand{\xClassCogMaxMcc}{0.234}
\newcommand{\xClassCogMaxRank}{6}
\newcommand{\xClassRepos}{205}
\newcommand{\xClassCandidates}{18}
\newcommand{\xMethod}{cognitive}
\newcommand{\xMethodMcc}{0.600}
\newcommand{\xMethodSecond}{cc}
\newcommand{\xMethodSecondMcc}{0.531}
\newcommand{\xMethodRepos}{212}
\newcommand{\gateRequired}{72}
\newcommand{\gateHalfWidth}{0.095}
\newcommand{\gateExploreEstimate}{$-$0.063 [$-$0.155, +0.034]}
\newcommand{\gateExploreRepos}{205}
\newcommand{\gateExplorePositiveRepos}{80}
\newcommand{\gateVOneRequired}{73}
\newcommand{\gateVOneS}{0.434}
\newcommand{\primRepos}{208}
\newcommand{\primReposWithPositive}{66}
\newcommand{\primUniv}{934}
\newcommand{\primPos}{98}
\newcommand{\primNeg}{836}
\newcommand{\primBaseRate}{10.5\%}
\newcommand{\primMicro}{$-$0.041 [$-$0.130, +0.054]}
\newcommand{\primMicroValue}{$-$0.041}
\newcommand{\primMicroLo}{$-$0.130}
\newcommand{\primMicroHi}{+0.054}
\newcommand{\primReading}{no difference resolved}
\newcommand{\primVerdict}{no difference resolved (underpowered)}
\newcommand{\primMacro}{$-$0.004 [$-$0.032, +0.026]}
\newcommand{\primHalfWidthConfirm}{0.092}
\newcommand{\primGateReposWithBlob}{66}
\newcommand{\primXMcc}{0.187}
\newcommand{\primXPrecision}{0.26}
\newcommand{\primXRecall}{0.31}
\newcommand{\primXFpr}{0.104}
\newcommand{\primXTp}{30}
\newcommand{\primXFp}{87}
\newcommand{\primXFn}{68}
\newcommand{\primXTn}{749}
\newcommand{\primXFlagged}{117}
\newcommand{\primDMcc}{0.146}
\newcommand{\primDPrecision}{0.19}
\newcommand{\primDRecall}{0.41}
\newcommand{\primDFpr}{0.208}
\newcommand{\primDTp}{40}
\newcommand{\primDFp}{174}
\newcommand{\primDFn}{58}
\newcommand{\primDTn}{662}
\newcommand{\primDFlagged}{214}
\newcommand{\primAucX}{0.768}
\newcommand{\primAucD}{0.631}
\newcommand{\primAucDiff}{$-$0.137 [$-$0.195, $-$0.080]}
\newcommand{\primJpMicro}{$-$0.050 [$-$0.142, +0.049]}
\newcommand{\primJpVerdict}{no difference resolved (underpowered)}
\newcommand{\primJpUniv}{933}
\newcommand{\primJpAucX}{0.767}
\newcommand{\primJpAucD}{0.631}
\newcommand{\primJpAucDiff}{$-$0.136 [$-$0.194, $-$0.079]}
\newcommand{\primCtorMicro}{$-$0.062 [$-$0.130, +0.011]}
\newcommand{\primCtorReading}{no difference resolved}
\newcommand{\primCtorAgrees}{agrees}
\newcommand{\primCtorMcc}{0.125}
\newcommand{\primCtorPrecision}{0.20}
\newcommand{\primCtorRecall}{0.27}
\newcommand{\primMethodsVsCtorMicro}{+0.021 [$-$0.048, +0.097]}
\newcommand{\primEMicro}{+0.025 [$-$0.052, +0.100]}
\newcommand{\primEReading}{no difference resolved}
\newcommand{\primEMacro}{+0.011 [$-$0.014, +0.038]}
\newcommand{\primEYMcc}{0.212}
\newcommand{\primEYPrecision}{0.26}
\newcommand{\primEYRecall}{0.37}
\newcommand{\primEYFp}{102}
\newcommand{\primEYTp}{36}
\newcommand{\primAmpMicro}{+0.009 [$-$0.075, +0.104]}
\newcommand{\primAmpReading}{no difference resolved}
\newcommand{\primAmpMacro}{+0.019 [$-$0.013, +0.056]}
\newcommand{\primAmpYMcc}{0.196}
\newcommand{\primAmpYPrecision}{0.20}
\newcommand{\primAmpYRecall}{0.56}
\newcommand{\primAmpYFp}{224}
\newcommand{\primAmpYTp}{55}
\newcommand{\primDvsEMicro}{$-$0.066 [$-$0.130, $-$0.000]}
\newcommand{\primDvsEReading}{exemption only beats the whole factor}
\newcommand{\primDvsEMacro}{$-$0.015 [$-$0.038, +0.008]}
\newcommand{\primDvsEYMcc}{0.146}
\newcommand{\primDvsEYPrecision}{0.19}
\newcommand{\primDvsEYRecall}{0.41}
\newcommand{\primDvsEYFp}{174}
\newcommand{\primDvsEYTp}{40}
\newcommand{\primDvsAMicro}{$-$0.051 [$-$0.112, +0.006]}
\newcommand{\primDvsAReading}{no difference resolved}
\newcommand{\primDvsAMacro}{$-$0.023 [$-$0.046, $-$0.004]}
\newcommand{\primDvsAYMcc}{0.146}
\newcommand{\primDvsAYPrecision}{0.19}
\newcommand{\primDvsAYRecall}{0.41}
\newcommand{\primDvsAYFp}{174}
\newcommand{\primDvsAYTp}{40}
\newcommand{\primNvsXMicro}{+0.049 [$-$0.023, +0.121]}
\newcommand{\primNvsXReading}{no difference resolved}
\newcommand{\primNvsXMacro}{+0.029 [0.000, +0.058]}
\newcommand{\primDvsNMicro}{$-$0.091 [$-$0.175, $-$0.006]}
\newcommand{\primDvsNReading}{the size factor beats the coupling factor}
\newcommand{\primDvsNMacro}{$-$0.033 [$-$0.072, +0.004]}
\newcommand{\primNMcc}{0.236}
\newcommand{\primNPrecision}{0.22}
\newcommand{\primNRecall}{0.60}
\newcommand{\primNTp}{59}
\newcommand{\primNFp}{211}
\newcommand{\primMechanism}{not read: the primary does not read refinement}
\newcommand{\primJpDvsAReading}{amplification only beats the whole factor}
\newcommand{\primJpDvsAMicro}{$-$0.062 [$-$0.129, $-$0.004]}
\newcommand{\primJpDvsEMicro}{$-$0.078 [$-$0.148, $-$0.008]}
\newcommand{\primJpDvsNMicro}{$-$0.097 [$-$0.184, $-$0.009]}
\newcommand{\primJpEMicro}{+0.028 [$-$0.048, +0.102]}
\newcommand{\primJpAmpMicro}{+0.012 [$-$0.074, +0.105]}
\newcommand{\leafCatchFlagged}{38}
\newcommand{\leafCatchPositive}{8}
\newcommand{\leafCatchRate}{21.1\%}
\newcommand{\leafRestFlagged}{79}
\newcommand{\leafRestPositive}{22}
\newcommand{\leafRestRate}{27.8\%}
\newcommand{\leafDiff}{$-$0.068 [$-$0.238, +0.102]}
\newcommand{\leafReading}{unresolved}
\newcommand{\leafMethodCatchFlagged}{28}
\newcommand{\leafMethodCatchPositive}{13}
\newcommand{\leafMethodCatchRate}{46.4\%}
\newcommand{\leafMethodRestRate}{76.8\%}
\newcommand{\leafMethodDiff}{$-$0.304 [$-$0.541, $-$0.082]}
\newcommand{\leafClassUnresolved}{12}
\newcommand{\leafClassVariants}{12}
\newcommand{\leafMethodReading}{supported}
\newcommand{\leafMethodSupported}{3}
\newcommand{\leafMethodVariants}{3}
\newcommand{\leafFontanaReading}{supported}
\newcommand{\leafFontanaSupported}{8}
\newcommand{\leafFontanaFlowSupported}{4}
\newcommand{\leafFontanaVariants}{12}
\newcommand{\leafMethodDataDiff}{$-$0.208 [$-$0.413, $-$0.017]}
\newcommand{\leafMethodMessageDiff}{$-$0.327 [$-$0.557, $-$0.102]}
\newcommand{\leafFontanaCatchRate}{75.0\%}
\newcommand{\leafFontanaRestRate}{94.9\%}
\newcommand{\leafFontanaDiff}{$-$0.199 [$-$0.362, $-$0.001]}
\newcommand{\leafFontanaCatchFlagged}{20}
\newcommand{\leafFontanaCatchPositive}{15}
\newcommand{\consVerdict}{outgoing data only}
\newcommand{\consCtorVerdict}{outgoing data only}
\newcommand{\consCtorAgrees}{agrees}
\newcommand{\consJpVerdict}{outgoing data only}
\newcommand{\consJpAgrees}{agrees}
\newcommand{\pOneHolds}{does not hold}
\newcommand{\pOneDPeak}{+0.041 [$-$0.031, +0.114]}
\newcommand{\pOneDDensity}{+0.094 [+0.013, +0.174]}
\newcommand{\pOneOPeak}{+0.125 [+0.054, +0.195]}
\newcommand{\pOneODensity}{+0.151 [+0.076, +0.227]}
\newcommand{\pOneOHolds}{holds}
\newcommand{\pTwoHolds}{holds}
\newcommand{\pTwoBlob}{+0.073 [+0.024, +0.122]}
\newcommand{\pTwoLong}{+0.188 [+0.109, +0.266]}
\newcommand{\pTwoFontana}{+0.242 [+0.195, +0.295]}
\newcommand{\pTwoMinToolLo}{+0.013}
\newcommand{\pTwoJpBlob}{+0.065 [+0.019, +0.111]}
\newcommand{\pTwoJpHolds}{holds}
\newcommand{\pThreeReads}{outgoing-only}
\newcommand{\pThreeSize}{$-$0.080 [$-$0.144, $-$0.014]}
\newcommand{\pThreePeak}{$-$0.084 [$-$0.139, $-$0.029]}
\newcommand{\pThreeDensity}{$-$0.058 [$-$0.110, $-$0.009]}
\newcommand{\pFourPositive}{5}
\newcommand{\pFourCells}{12}
\newcommand{\pFourMajority}{not a majority}
\newcommand{\pFourLoAbove}{5}
\newcommand{\pFourHiBelow}{4}
\newcommand{\pOneDataHolds}{does not hold}
\newcommand{\pOneDataOHolds}{holds}
\newcommand{\pThreeDataReads}{outgoing-only}
\newcommand{\pOneDataDPeak}{+0.061 [$-$0.014, +0.137]}
\newcommand{\pOneDataDDensity}{+0.109 [+0.024, +0.192]}
\newcommand{\pOneDataOPeak}{+0.103 [+0.033, +0.176]}
\newcommand{\pOneDataODensity}{+0.131 [+0.054, +0.204]}
\newcommand{\pOneMessageHolds}{does not hold}
\newcommand{\pOneMessageOHolds}{holds}
\newcommand{\pThreeMessageReads}{outgoing-only}
\newcommand{\pOneMessageDPeak}{+0.066 [$-$0.001, +0.145]}
\newcommand{\pOneMessageDDensity}{+0.107 [+0.026, +0.189]}
\newcommand{\pOneMessageOPeak}{+0.130 [+0.059, +0.199]}
\newcommand{\pOneMessageODensity}{+0.147 [+0.072, +0.219]}
\newcommand{\cellBlobSizeD}{$-$0.032 [$-$0.114, +0.048]}
\newcommand{\cellBlobSizeO}{+0.048 [$-$0.014, +0.108]}
\newcommand{\cellBlobSizeDO}{$-$0.080 [$-$0.144, $-$0.014]}
\newcommand{\cellBlobSizeN}{208}
\newcommand{\cellBlobPeakD}{+0.041 [$-$0.031, +0.114]}
\newcommand{\cellBlobPeakO}{+0.125 [+0.054, +0.195]}
\newcommand{\cellBlobPeakDO}{$-$0.084 [$-$0.139, $-$0.029]}
\newcommand{\cellBlobPeakN}{208}
\newcommand{\cellBlobDensityD}{+0.094 [+0.013, +0.174]}
\newcommand{\cellBlobDensityO}{+0.151 [+0.076, +0.227]}
\newcommand{\cellBlobDensityDO}{$-$0.058 [$-$0.110, $-$0.009]}
\newcommand{\cellBlobDensityN}{208}
\newcommand{\cellLongSizeD}{$-$0.232 [$-$0.317, $-$0.144]}
\newcommand{\cellLongSizeO}{$-$0.159 [$-$0.200, $-$0.121]}
\newcommand{\cellLongSizeDO}{$-$0.073 [$-$0.156, +0.012]}
\newcommand{\cellLongSizeN}{202}
\newcommand{\cellLongDensityD}{+0.280 [+0.197, +0.362]}
\newcommand{\cellLongDensityO}{+0.222 [+0.162, +0.280]}
\newcommand{\cellLongDensityDO}{+0.058 [+0.006, +0.110]}
\newcommand{\cellLongDensityN}{202}
\newcommand{\cellLongStructureD}{$-$0.044 [$-$0.115, +0.028]}
\newcommand{\cellLongStructureO}{$-$0.028 [$-$0.086, +0.034]}
\newcommand{\cellLongStructureDO}{$-$0.017 [$-$0.061, +0.024]}
\newcommand{\cellLongStructureN}{202}
\newcommand{\cellFontanaSizeD}{$-$0.073 [$-$0.193, +0.037]}
\newcommand{\cellFontanaSizeO}{+0.063 [$-$0.004, +0.121]}
\newcommand{\cellFontanaSizeDO}{$-$0.136 [$-$0.243, $-$0.038]}
\newcommand{\cellFontanaSizeN}{72}
\newcommand{\cellFontanaPeakD}{+0.169 [+0.068, +0.269]}
\newcommand{\cellFontanaPeakO}{+0.315 [+0.249, +0.381]}
\newcommand{\cellFontanaPeakDO}{$-$0.146 [$-$0.230, $-$0.067]}
\newcommand{\cellFontanaPeakN}{72}
\newcommand{\cellFontanaDensityD}{+0.413 [+0.326, +0.503]}
\newcommand{\cellFontanaDensityO}{+0.573 [+0.499, +0.658]}
\newcommand{\cellFontanaDensityDO}{$-$0.160 [$-$0.247, $-$0.079]}
\newcommand{\cellFontanaDensityN}{72}
\newcommand{\cellsCountedFlow}{21}
\newcommand{\cellsCountedDLoAbove}{9}
\newcommand{\cellsCountedDHiBelow}{5}
\newcommand{\cellsCountedOLoAbove}{12}
\newcommand{\cellsCountedOHiBelow}{6}
\newcommand{\cellsCountedDOHiBelow}{12}
\newcommand{\cellsCountedDOLoAbove}{5}
\newcommand{\cellsCountedClassFlow}{12}
\newcommand{\cellsCountedClassDOHiBelow}{12}
\newcommand{\cellsDensityFlow}{7}
\newcommand{\cellsDensityDLoAbove}{7}
\newcommand{\cellsSizeFlow}{7}
\newcommand{\cellsSizeDHiBelow}{5}
\newcommand{\betweenCells}{34}
\newcommand{\betweenFlowDataLoAbove}{10}
\newcommand{\betweenFlowDataHiBelow}{1}
\newcommand{\betweenMessageFlowLoAbove}{2}
\newcommand{\betweenMessageFlowHiBelow}{0}
\newcommand{\betweenBlobPeakMessageFlow}{+0.025 [+0.004, +0.050]}
\newcommand{\betweenLongStructureMessageFlow}{+0.019 [$-$0.004, +0.042]}
\newcommand{\betweenFontanaDensityMessageFlow}{+0.045 [+0.003, +0.100]}
\newcommand{\betweenBlobPeakFlowData}{$-$0.020 [$-$0.054, +0.014]}
\newcommand{\secDDataMicro}{$-$0.008 [$-$0.101, +0.096]}
\newcommand{\secDDataMacro}{+0.013 [$-$0.018, +0.048]}
\newcommand{\secDDataReading}{no difference resolved}
\newcommand{\secDDataMcc}{0.180}
\newcommand{\secDDataPrecision}{0.20}
\newcommand{\secDDataRecall}{0.47}
\newcommand{\secDDataFp}{182}
\newcommand{\secDDataTp}{46}
\newcommand{\secDMessageMicro}{$-$0.005 [$-$0.089, +0.084]}
\newcommand{\secDMessageMacro}{+0.012 [$-$0.017, +0.044]}
\newcommand{\secDMessageReading}{no difference resolved}
\newcommand{\secDMessageMcc}{0.182}
\newcommand{\secDMessagePrecision}{0.20}
\newcommand{\secDMessageRecall}{0.48}
\newcommand{\secDMessageFp}{186}
\newcommand{\secDMessageTp}{47}
\newcommand{\secOFlowMicro}{+0.085 [+0.001, +0.182]}
\newcommand{\secOFlowMacro}{+0.063 [+0.028, +0.103]}
\newcommand{\secOFlowReading}{the divergence factor refines complexity}
\newcommand{\secOFlowMcc}{0.272}
\newcommand{\secOFlowPrecision}{0.23}
\newcommand{\secOFlowRecall}{0.69}
\newcommand{\secOFlowFp}{233}
\newcommand{\secOFlowTp}{68}
\newcommand{\secODataMicro}{+0.035 [$-$0.053, +0.134]}
\newcommand{\secODataMacro}{+0.040 [+0.006, +0.076]}
\newcommand{\secODataReading}{no difference resolved}
\newcommand{\secODataMcc}{0.222}
\newcommand{\secODataPrecision}{0.21}
\newcommand{\secODataRecall}{0.61}
\newcommand{\secODataFp}{231}
\newcommand{\secODataTp}{60}
\newcommand{\secOMessageMicro}{+0.087 [+0.003, +0.183]}
\newcommand{\secOMessageMacro}{+0.058 [+0.025, +0.097]}
\newcommand{\secOMessageReading}{the divergence factor refines complexity}
\newcommand{\secOMessageMcc}{0.274}
\newcommand{\secOMessagePrecision}{0.23}
\newcommand{\secOMessageRecall}{0.68}
\newcommand{\secOMessageFp}{225}
\newcommand{\secOMessageTp}{67}
\newcommand{\secDAloneMicro}{$-$0.014 [$-$0.109, +0.091]}
\newcommand{\secDAloneMacro}{+0.004 [$-$0.027, +0.038]}
\newcommand{\secDAloneReading}{no difference resolved}
\newcommand{\secDAloneMcc}{0.173}
\newcommand{\secDAlonePrecision}{0.21}
\newcommand{\secDAloneRecall}{0.40}
\newcommand{\secDAloneFp}{145}
\newcommand{\secDAloneTp}{39}
\newcommand{\secOAloneMicro}{0.000 [$-$0.100, +0.106]}
\newcommand{\secOAloneMacro}{+0.019 [$-$0.017, +0.058]}
\newcommand{\secOAloneReading}{no difference resolved}
\newcommand{\secOAloneMcc}{0.187}
\newcommand{\secOAlonePrecision}{0.21}
\newcommand{\secOAloneRecall}{0.45}
\newcommand{\secOAloneFp}{163}
\newcommand{\secOAloneTp}{44}
\newcommand{\secAndMicro}{$-$0.023 [$-$0.092, +0.045]}
\newcommand{\secAndMacro}{$-$0.015 [$-$0.037, +0.004]}
\newcommand{\secAndReading}{no difference resolved}
\newcommand{\secAndMcc}{0.164}
\newcommand{\secAndPrecision}{0.30}
\newcommand{\secAndRecall}{0.17}
\newcommand{\secAndFp}{39}
\newcommand{\secAndTp}{17}
\newcommand{\secDPubMicro}{$-$0.027 [$-$0.116, +0.071]}
\newcommand{\secDPubMacro}{+0.003 [$-$0.025, +0.034]}
\newcommand{\secDPubReading}{no difference resolved}
\newcommand{\secDPubMcc}{0.161}
\newcommand{\secDPubPrecision}{0.18}
\newcommand{\secDPubRecall}{0.49}
\newcommand{\secDPubFp}{213}
\newcommand{\secDPubTp}{48}
\newcommand{\secDFormBMicro}{+0.044 [$-$0.038, +0.139]}
\newcommand{\secDFormBMacro}{+0.035 [+0.003, +0.073]}
\newcommand{\secDFormBReading}{no difference resolved}
\newcommand{\secDFormBMcc}{0.231}
\newcommand{\secDFormBPrecision}{0.21}
\newcommand{\secDFormBRecall}{0.60}
\newcommand{\secDFormBFp}{216}
\newcommand{\secDFormBTp}{59}
\newcommand{\secDFormBCtorMicro}{+0.064 [$-$0.023, +0.162]}
\newcommand{\secDFormBCtorMacro}{+0.047 [+0.014, +0.087]}
\newcommand{\secDFormBCtorReading}{no difference resolved}
\newcommand{\secDFormBCtorMcc}{0.251}
\newcommand{\secDFormBCtorPrecision}{0.22}
\newcommand{\secDFormBCtorRecall}{0.63}
\newcommand{\secDFormBCtorFp}{216}
\newcommand{\secDFormBCtorTp}{62}
\newcommand{\secOFormBMicro}{+0.026 [$-$0.046, +0.107]}
\newcommand{\secOFormBMacro}{+0.027 [$-$0.004, +0.061]}
\newcommand{\secOFormBReading}{no difference resolved}
\newcommand{\secOFormBMcc}{0.213}
\newcommand{\secOFormBPrecision}{0.21}
\newcommand{\secOFormBRecall}{0.54}
\newcommand{\secOFormBFp}{195}
\newcommand{\secOFormBTp}{53}
\newcommand{\secOtherArms}{41}
\newcommand{\secOtherLoAbove}{2}
\newcommand{\secOtherHiBelow}{0}
\newcommand{\secOtherMacroLoAbove}{12}
\newcommand{\secFormBArms}{15}
\newcommand{\secFormBMacroLoAbove}{9}
\newcommand{\secFormBMinMicro}{$-$0.003}
\newcommand{\secFormBMaxMicro}{+0.064}
\newcommand{\rhoXD}{0.21}
\newcommand{\rhoXOut}{0.55}
\newcommand{\rhoXDData}{0.10}
\newcommand{\rhoNomD}{$-$0.23}
\newcommand{\rhoNomOut}{0.40}
\newcommand{\rhoNlocSumOut}{0.57}
\newcommand{\rhoXDSeed}{0.16}
\newcommand{\rhoXOutSeed}{0.52}
\newcommand{\rhoXmD}{0.25}
\newcommand{\rhoXmOut}{0.39}
\newcommand{\trackClosely}{do not track closely}
\newcommand{\tertileOne}{$-$0.012 [$-$0.021, 0.000]}
\newcommand{\tertileTwo}{$-$0.026 [$-$0.115, +0.085]}
\newcommand{\tertileThree}{$-$0.022 [$-$0.151, +0.118]}
\newcommand{\methRepos}{202}
\newcommand{\methUniv}{1033}
\newcommand{\methPos}{125}
\newcommand{\methBaseRate}{12.1\%}
\newcommand{\methXMcc}{0.603}
\newcommand{\methXPrecision}{0.69}
\newcommand{\methXRecall}{0.61}
\newcommand{\methXTp}{76}
\newcommand{\methXFp}{34}
\newcommand{\methXFn}{49}
\newcommand{\methDMcc}{0.511}
\newcommand{\methDPrecision}{0.86}
\newcommand{\methDRecall}{0.34}
\newcommand{\methDTp}{43}
\newcommand{\methDFp}{7}
\newcommand{\methDFn}{82}
\newcommand{\methMicro}{$-$0.092 [$-$0.188, +0.010]}
\newcommand{\methReading}{no difference resolved}
\newcommand{\methMacro}{$-$0.057 [$-$0.099, $-$0.015]}
\newcommand{\methNoHelpMicro}{$-$0.144 [$-$0.244, $-$0.043]}
\newcommand{\methNoHelpReading}{it makes it worse}
\newcommand{\methNoHelpMcc}{0.460}
\newcommand{\methNoHelpVsAllMicro}{$-$0.051 [$-$0.125, +0.016]}
\newcommand{\methEMicro}{$-$0.002 [$-$0.077, +0.073]}
\newcommand{\methEReading}{no difference resolved}
\newcommand{\methEMcc}{0.601}
\newcommand{\methERecall}{0.58}
\newcommand{\methAmpMicro}{$-$0.093 [$-$0.174, $-$0.007]}
\newcommand{\methAmpReading}{the amplification makes it worse}
\newcommand{\methAmpMcc}{0.510}
\newcommand{\methAmpRecall}{0.40}
\newcommand{\methDvsEMicro}{$-$0.090 [$-$0.173, $-$0.006]}
\newcommand{\methDvsEReading}{exemption only beats the whole factor}
\newcommand{\methDvsAMicro}{+0.001 [$-$0.043, +0.040]}
\newcommand{\methDvsAReading}{no difference resolved}
\newcommand{\methNeVsXMicro}{$-$0.060 [$-$0.153, +0.045]}
\newcommand{\methNeVsXReading}{no difference resolved}
\newcommand{\methNeVsAllMicro}{+0.032 [$-$0.058, +0.126]}
\newcommand{\methNeVsAllMacro}{+0.037 [+0.004, +0.071]}
\newcommand{\methNeMcc}{0.543}
\newcommand{\methNePrecision}{0.54}
\newcommand{\methNeRecall}{0.68}
\newcommand{\methMechanism}{not read: the primary does not read refinement}
\newcommand{\methJpMicro}{$-$0.079 [$-$0.171, +0.026]}
\newcommand{\methJpReading}{no difference resolved}
\newcommand{\toolPmdMicro}{$-$0.143 [$-$0.176, $-$0.106]}
\newcommand{\toolPmdReading}{it makes it worse}
\newcommand{\toolPmdN}{490}
\newcommand{\toolPmdXMcc}{0.351}
\newcommand{\toolPmdDMcc}{0.208}
\newcommand{\toolPmdXRecall}{0.43}
\newcommand{\toolPmdDRecall}{0.52}
\newcommand{\toolPmdXPrecision}{0.37}
\newcommand{\toolPmdDPrecision}{0.17}
\newcommand{\toolPmdCtorMicro}{$-$0.117 [$-$0.153, $-$0.078]}
\newcommand{\toolDecorMicro}{$-$0.093 [$-$0.121, $-$0.059]}
\newcommand{\toolDecorReading}{it makes it worse}
\newcommand{\toolDecorN}{490}
\newcommand{\toolDecorXMcc}{0.202}
\newcommand{\toolDecorDMcc}{0.109}
\newcommand{\toolDecorXRecall}{0.25}
\newcommand{\toolDecorDRecall}{0.35}
\newcommand{\toolDecorXPrecision}{0.29}
\newcommand{\toolDecorDPrecision}{0.15}
\newcommand{\toolDecorCtorMicro}{$-$0.051 [$-$0.097, $-$0.002]}
\newcommand{\toolWmcMicro}{$-$0.144 [$-$0.185, $-$0.099]}
\newcommand{\toolWmcReading}{it makes it worse}
\newcommand{\toolWmcN}{491}
\newcommand{\toolWmcXMcc}{0.390}
\newcommand{\toolWmcDMcc}{0.247}
\newcommand{\toolWmcXRecall}{0.38}
\newcommand{\toolWmcDRecall}{0.50}
\newcommand{\toolWmcXPrecision}{0.52}
\newcommand{\toolWmcDPrecision}{0.25}
\newcommand{\toolWmcCtorMicro}{$-$0.122 [$-$0.161, $-$0.074]}
\newcommand{\toolNpathMicro}{$-$0.021 [$-$0.054, +0.013]}
\newcommand{\toolNpathReading}{no difference resolved}
\newcommand{\toolNpathN}{490}
\newcommand{\toolNpathXMcc}{0.350}
\newcommand{\toolNpathDMcc}{0.330}
\newcommand{\toolNpathXRecall}{0.97}
\newcommand{\toolNpathDRecall}{0.62}
\newcommand{\toolNpathXPrecision}{0.13}
\newcommand{\toolNpathDPrecision}{0.19}
\newcommand{\toolBrainMicro}{+0.036 [+0.025, +0.047]}
\newcommand{\toolBrainReading}{the divergence factor refines complexity}
\newcommand{\toolBrainN}{491}
\newcommand{\toolBrainXMcc}{0.220}
\newcommand{\toolBrainDMcc}{0.256}
\newcommand{\toolBrainXRecall}{1.00}
\newcommand{\toolBrainDRecall}{0.78}
\newcommand{\toolBrainXPrecision}{0.05}
\newcommand{\toolBrainDPrecision}{0.09}
\newcommand{\toolBrainJpMicro}{+0.042 [+0.032, +0.053]}
\newcommand{\fontanaPopSystems}{72}
\newcommand{\fontanaPopUniv}{380}
\newcommand{\fontanaPopPos}{139}
\newcommand{\fontanaMicro}{$-$0.093 [$-$0.220, +0.035]}
\newcommand{\fontanaReading}{no difference resolved}
\newcommand{\fontanaMacro}{$-$0.065 [$-$0.155, +0.021]}
\newcommand{\fontanaXMcc}{0.765}
\newcommand{\fontanaDMcc}{0.672}
\newcommand{\fontanaXPrecision}{0.91}
\newcommand{\fontanaDPrecision}{0.79}
\newcommand{\fontanaXRecall}{0.78}
\newcommand{\fontanaDRecall}{0.80}
\newcommand{\fontanaXFp}{11}
\newcommand{\fontanaDFp}{30}
\newcommand{\fontanaCtorMicro}{$-$0.065 [$-$0.169, +0.039]}
\newcommand{\fontanaCtorReading}{no difference resolved}
\newcommand{\fontanaCtorMacro}{$-$0.084 [$-$0.167, $-$0.004]}
\newcommand{\fontanaVOneSystems}{19}
\newcommand{\fontanaVOneUniv}{100}
\newcommand{\fontanaVOnePos}{46}
\newcommand{\fontanaVOneMicro}{$-$0.203 [$-$0.337, $-$0.074]}
\newcommand{\fontanaVOneReading}{it makes it worse}
\newcommand{\fontanaVOneMeasured}{46}
\newcommand{\primVOneMicro}{$-$0.075 [$-$0.191, +0.045]}
\newcommand{\consVOneVerdict}{outgoing data only}
\newcommand{\fontanaPmdMedianMeasurer}{0.978}
\newcommand{\cellFontanaDensityOFlag}{+0.573 [+0.499, +0.658]}
\newcommand{\friedBlobSizeRankX}{2.08}
\newcommand{\friedBlobSizeRankD}{2.16}
\newcommand{\friedBlobSizeRankO}{1.76}
\newcommand{\friedBlobSizeP}{$p = 0.100$}
\newcommand{\friedBlobSizeN}{53}
\newcommand{\friedBlobPeakRankX}{2.32}
\newcommand{\friedBlobPeakRankD}{2.11}
\newcommand{\friedBlobPeakRankO}{1.58}
\newcommand{\friedBlobPeakP}{$p < 0.001$}
\newcommand{\friedBlobPeakN}{52}
\newcommand{\friedBlobDensityRankX}{2.37}
\newcommand{\friedBlobDensityRankD}{2.00}
\newcommand{\friedBlobDensityRankO}{1.63}
\newcommand{\friedBlobDensityP}{$p < 0.001$}
\newcommand{\friedBlobDensityN}{53}
\newcommand{\friedLongSizeRankX}{1.40}
\newcommand{\friedLongSizeRankD}{2.40}
\newcommand{\friedLongSizeRankO}{2.20}
\newcommand{\friedLongSizeP}{$p < 0.001$}
\newcommand{\friedLongSizeN}{63}
\newcommand{\friedLongDensityRankX}{2.55}
\newcommand{\friedLongDensityRankD}{1.71}
\newcommand{\friedLongDensityRankO}{1.74}
\newcommand{\friedLongDensityP}{$p < 0.001$}
\newcommand{\friedLongDensityN}{66}
\newcommand{\friedLongStructureRankX}{1.95}
\newcommand{\friedLongStructureRankD}{2.14}
\newcommand{\friedLongStructureRankO}{1.92}
\newcommand{\friedLongStructureP}{$p = 0.387$}
\newcommand{\friedLongStructureN}{65}
\newcommand{\friedFontanaSizeRankX}{1.99}
\newcommand{\friedFontanaSizeRankD}{2.09}
\newcommand{\friedFontanaSizeRankO}{1.92}
\newcommand{\friedFontanaSizeP}{$p = 0.723$}
\newcommand{\friedFontanaSizeN}{44}
\newcommand{\friedFontanaPeakRankX}{2.65}
\newcommand{\friedFontanaPeakRankD}{1.85}
\newcommand{\friedFontanaPeakRankO}{1.50}
\newcommand{\friedFontanaPeakP}{$p < 0.001$}
\newcommand{\friedFontanaPeakN}{48}
\newcommand{\friedFontanaDensityRankX}{2.84}
\newcommand{\friedFontanaDensityRankD}{1.85}
\newcommand{\friedFontanaDensityRankO}{1.30}
\newcommand{\friedFontanaDensityP}{$p < 0.001$}
\newcommand{\friedFontanaDensityN}{48}
\newcommand{\matrixBlobRows}{16{,}450}
\newcommand{\matrixLongRows}{2{,}926}
\newcommand{\sensDiMicro}{$-$0.108 [$-$0.244, +0.036]}
\newcommand{\sensDiReading}{no difference resolved}
\newcommand{\sensDiRepos}{116}
\newcommand{\sensDiXMcc}{0.225}
\newcommand{\sensDiDMcc}{0.117}
\newcommand{\sensDiUniv}{441}
\newcommand{\sensNeVsXMicro}{+0.017 [$-$0.072, +0.117]}
\newcommand{\sensNeVsXReading}{no difference resolved}
\newcommand{\sensNeVsAllMicro}{+0.058 [+0.001, +0.117]}
\newcommand{\sensNeVsAllReading}{fitting on the non-exempt units is better}
\newcommand{\sensNeMcc}{0.204}
\newcommand{\sensNePrecision}{0.25}
\newcommand{\sensNeRecall}{0.37}
\newcommand{\sensNeFp}{107}
\newcommand{\sensJpNeVsAllMicro}{+0.068 [+0.006, +0.134]}
\newcommand{\sensJpNeVsAllReading}{fitting on the non-exempt units is better}
\newcommand{\sensUncoupledMicro}{$-$0.041 [$-$0.130, +0.054]}
\newcommand{\sensBudgetRepos}{0}
\newcommand{\interBlobBothTp}{17}
\newcommand{\interBlobBothFp}{39}
\newcommand{\interBlobBothPrecision}{30.4\%}
\newcommand{\interBlobXOnlyPrecision}{21.3\%}
\newcommand{\interBlobDOnlyPrecision}{21.7\%}
\newcommand{\interBlobEitherMcc}{0.223}
\newcommand{\interBlobBothVsX}{$-$0.023 [$-$0.092, +0.045]}
\newcommand{\interBlobEitherVsX}{+0.036 [$-$0.023, +0.106]}
\newcommand{\interBlobDVsX}{+0.007 [$-$0.080, +0.104]}
\newcommand{\interBlobDMcc}{0.194}
\newcommand{\interLongBothTp}{47}
\newcommand{\interLongBothFp}{7}
\newcommand{\interLongBothPrecision}{87.0\%}
\newcommand{\interLongDOnlyPrecision}{27.3\%}
\newcommand{\interLongBothVsX}{$-$0.064 [$-$0.147, +0.015]}
\newcommand{\interLongEitherVsX}{$-$0.028 [$-$0.091, +0.043]}
\newcommand{\interLongDVsX}{$-$0.139 [$-$0.238, $-$0.032]}
\newcommand{\interLongDMcc}{0.464}
\newcommand{\interFontanaEitherVsX}{+0.052 [+0.011, +0.094]}
\newcommand{\interFontanaEitherMcc}{0.818}
\newcommand{\interFontanaBothVsX}{$-$0.066 [$-$0.126, $-$0.005]}
\newcommand{\interFontanaBothPrecision}{98.8\%}
\newcommand{\primJpReposWithPositive}{66}
\newcommand{\analysisVTwoShaShort}{9519e2466f44}
\newcommand{\priorMadeUtc}{2026-10-26 09:49:15 UTC}
\newcommand{\priorPrimRepos}{67}
\newcommand{\priorPrimUniv}{548}
\newcommand{\priorPrimPos}{98}
\newcommand{\priorPrimMicro}{$-$0.075 [$-$0.191, +0.045]}
\newcommand{\priorPrimVerdict}{no difference resolved (underpowered)}
\newcommand{\priorPrimXMcc}{0.212}
\newcommand{\priorPrimDMcc}{0.136}
\newcommand{\priorPrimXPrecision}{0.38}
\newcommand{\priorPrimDPrecision}{0.26}
\newcommand{\priorPrimAucX}{0.767}
\newcommand{\priorPrimAucD}{0.620}
\newcommand{\priorPrimAucDiff}{$-$0.147 [$-$0.201, $-$0.093]}
\newcommand{\priorDvsNMicro}{$-$0.127 [$-$0.230, $-$0.021]}
\newcommand{\priorDvsEMicro}{$-$0.087 [$-$0.166, $-$0.012]}
\newcommand{\priorGateReposWithBlob}{66}
\newcommand{\priorGateRequired}{108}
\newcommand{\priorLeafDiff}{$-$0.016 [$-$0.273, +0.271]}
\newcommand{\priorLeafCatchRate}{36.4\%}
\newcommand{\priorLeafRestRate}{37.9\%}
\newcommand{\priorConsVerdict}{outgoing data only}
\newcommand{\priorPOneDPeak}{+0.018 [$-$0.068, +0.112]}
\newcommand{\priorPOneDDensity}{+0.062 [$-$0.041, +0.161]}
\newcommand{\priorPOneOPeak}{+0.147 [+0.061, +0.234]}
\newcommand{\priorPOneODensity}{+0.162 [+0.068, +0.255]}
\newcommand{\priorPThreePeak}{$-$0.128 [$-$0.199, $-$0.062]}
\newcommand{\priorMethRepos}{78}
\newcommand{\priorMethUniv}{711}
\newcommand{\priorMethPos}{125}
\newcommand{\priorMethMicro}{$-$0.136 [$-$0.231, $-$0.041]}
\newcommand{\priorMethReading}{it makes it worse}
\newcommand{\priorFontanaSystems}{51}
\newcommand{\priorFontanaUniv}{305}
\newcommand{\priorFontanaMicro}{$-$0.118 [$-$0.260, +0.020]}
\newcommand{\priorFontanaReading}{no difference resolved}
\newcommand{\priorJpPrimMicro}{$-$0.084 [$-$0.205, +0.032]}
\newcommand{\restoredBlobRepos}{141}
\newcommand{\restoredBlobUnits}{386}
\newcommand{\restoredLongRepos}{124}
\newcommand{\restoredLongUnits}{322}
\newcommand{\restoredFontanaSystems}{21}
\newcommand{\restoredFontanaUnits}{75}
\newcommand{\priorXClass}{nloc\_max}
\newcommand{\priorXClassMcc}{0.277}
\newcommand{\priorXClassRepos}{81}
\newcommand{\priorXMethod}{cognitive}
\newcommand{\priorXMethodMcc}{0.638}
\newcommand{\priorXMethodRepos}{69}
\newcommand{\priorGateRequiredExplore}{108}
\newcommand{\priorGateHalfWidth}{0.116}
\newcommand{\priorGateExploreEstimate}{$-$0.063 [$-$0.178, +0.054]}
\newcommand{\priorGateExploreRepos}{81}
\newcommand{\priorDecisionsMadeUtc}{2026-10-26 07:54:14 UTC}
\newcommand{\restoredExploreBlobRepos}{124}
\newcommand{\restoredExploreLongRepos}{143}
\newcommand{\xClassChanged}{unchanged}
\newcommand{\xMethodChanged}{unchanged}
\newcommand{\expMadeUtc}{20261026 232945 UTC}
\newcommand{\expShaShort}{4bd479a1814c}
\newcommand{\expClassForms}{24}
\newcommand{\expMethodForms}{17}
\newcommand{\expBlobConfirmRepos}{208}
\newcommand{\expBlobExploreRepos}{205}
\newcommand{\expLongConfirmRepos}{202}
\newcommand{\expLongExploreRepos}{212}
\newcommand{\expClassTopConfirm}{lines summed $\times$ (1 + leaving calls)}
\newcommand{\expClassTopConfirmMcc}{0.286}
\newcommand{\expClassTopExplore}{lines summed $\times$ (1 + all internal calls)}
\newcommand{\expClassTopExploreMcc}{0.325}
\newcommand{\expSizeCallsConfirmMcc}{0.286}
\newcommand{\expSizeCallsExploreMcc}{0.314}
\newcommand{\expSizeCallsFontanaMcc}{0.822}
\newcommand{\expSizeCallsConfirmVsComplexity}{+0.032 [$-$0.037, +0.102]}
\newcommand{\expSizeCallsExploreVsComplexity}{+0.027 [$-$0.032, +0.086]}
\newcommand{\expSizeCallsConfirmVsSize}{+0.069 [+0.010, +0.135]}
\newcommand{\expSizeCallsExploreVsSize}{+0.028 [$-$0.021, +0.080]}
\newcommand{\expSizeCallsConfirmPrecision}{0.22}
\newcommand{\expSizeCallsConfirmRecall}{0.76}
\newcommand{\expSizeFanConfirmMcc}{0.280}
\newcommand{\expSizeFanExploreMcc}{0.317}
\newcommand{\expComplexityConfirmMcc}{0.254}
\newcommand{\expComplexityExploreMcc}{0.287}
\newcommand{\expSizeOnlyConfirmMcc}{0.217}
\newcommand{\expSizeOnlyExploreMcc}{0.286}
\newcommand{\expCallsAloneConfirmMcc}{0.214}
\newcommand{\expFanAloneConfirmMcc}{0.158}
\newcommand{\expDensityCallsConfirmMcc}{0.229}
\newcommand{\expDensityConfirmMcc}{0.039}
\newcommand{\expFormBSumConfirmMcc}{0.241}
\newcommand{\expRateConfirmMcc}{0.193}
\newcommand{\expXOConfirmVsSizeMax}{+0.036 [$-$0.031, +0.110]}
\newcommand{\expXOExploreVsSizeMax}{+0.016 [$-$0.050, +0.078]}
\newcommand{\expXOFontanaVsSizeMax}{+0.061 [+0.011, +0.110]}
\newcommand{\expFontanaTop}{fan-out (callee files) alone}
\newcommand{\expFontanaTopMcc}{0.847}
\newcommand{\expFontanaComplexityMcc}{0.835}
\newcommand{\expGateAnyConfirmMcc}{0.744}
\newcommand{\expGateAnyExploreMcc}{0.679}
\newcommand{\expGateAnyConfirmPrecision}{0.69}
\newcommand{\expGateAnyConfirmRecall}{0.87}
\newcommand{\expGateAnyConfirmVsTokens}{+0.119 [+0.070, +0.171]}
\newcommand{\expGateAnyExploreVsTokens}{+0.059 [$-$0.002, +0.118]}
\newcommand{\expGateAnyConfirmVsLines}{+0.224 [+0.164, +0.280]}
\newcommand{\expGateIntConfirmMcc}{0.697}
\newcommand{\expGateIntExploreMcc}{0.679}
\newcommand{\expGateIntConfirmVsTokens}{+0.072 [+0.009, +0.142]}
\newcommand{\expGateIntExploreVsTokens}{+0.059 [+0.004, +0.113]}
\newcommand{\expTokensConfirmMcc}{0.626}
\newcommand{\expTokensExploreMcc}{0.620}
\newcommand{\expLinesConfirmMcc}{0.520}
\newcommand{\expWeightConfirmMcc}{0.465}
\newcommand{\expWeightConfirmVsTokens}{$-$0.161 [$-$0.215, $-$0.110]}
\newcommand{\expDensityGateConfirmMcc}{0.017}
\newcommand{\expMethodDensityConfirmMcc}{$-$0.042}
\newcommand{\expGateLeaveConfirmMcc}{0.561}
\newcommand{\expToolPmdTop}{cyclomatic complexity, summed (WMC)}
\newcommand{\expToolPmdTopMcc}{0.582}
\newcommand{\expToolWmcTop}{cyclomatic complexity, summed (WMC)}
\newcommand{\expToolWmcTopMcc}{0.716}
\newcommand{\expToolDecorTop}{lines summed $\times$ number of methods (size only)}
\newcommand{\expToolDecorTopMcc}{0.455}
\newcommand{\expToolNpathCcMcc}{0.358}
\newcommand{\expToolNpathCcGateMcc}{0.339}
\newcommand{\expToolBrainCogMcc}{0.220}
\newcommand{\expToolBrainCogGateMcc}{0.210}
\newcommand{\matchedMadeUtc}{20261026 234905 UTC}
\newcommand{\matchedShaShort}{d6b77463f469}
\newcommand{\matchedClassFlags}{43{,}032}
\newcommand{\matchedClassUnits}{171{,}748}
\newcommand{\matchedClassRate}{25.1\%}
\newcommand{\matchedXAucConfirm}{0.768}
\newcommand{\matchedXAucExplore}{0.804}
\newcommand{\matchedOAucConfirm}{0.763}
\newcommand{\matchedOAucExplore}{0.774}
\newcommand{\matchedOAucVsXConfirm}{$-$0.005 [$-$0.043, +0.033]}
\newcommand{\matchedOAucVsXExplore}{$-$0.030 [$-$0.055, $-$0.006]}
\newcommand{\matchedNAucVsXConfirm}{+0.002 [$-$0.027, +0.035]}
\newcommand{\matchedDAucVsXConfirm}{$-$0.137 [$-$0.195, $-$0.078]}
\newcommand{\matchedSizeCallsAucVsXConfirm}{+0.008 [$-$0.028, +0.048]}
\newcommand{\matchedSizeCallsAucVsXExplore}{$-$0.011 [$-$0.038, +0.014]}
\newcommand{\matchedXOwnFlags}{117}
\newcommand{\matchedOOwnFlags}{301}
\newcommand{\matchedXOwnMcc}{0.187}
\newcommand{\matchedXMatchedMcc}{0.240}
\newcommand{\matchedXMatchedRecall}{0.64}
\newcommand{\matchedXMatchedPrecision}{0.21}
\newcommand{\matchedOMcc}{0.272}
\newcommand{\matchedOVsXConfirm}{+0.032 [$-$0.036, +0.099]}
\newcommand{\matchedOVsXExplore}{$-$0.015 [$-$0.067, +0.035]}
\newcommand{\matchedNVsXConfirm}{+0.052 [$-$0.005, +0.110]}
\newcommand{\matchedSizeCallsVsXConfirm}{+0.046 [$-$0.018, +0.111]}
\newcommand{\matchedSizeCallsVsXExplore}{$-$0.002 [$-$0.060, +0.055]}
\newcommand{\matchedDVsXConfirm}{$-$0.078 [$-$0.150, $-$0.003]}
\newcommand{\matchedFontanaOAucVsX}{+0.021 [+0.004, +0.043]}
\newcommand{\matchedFontanaOVsX}{+0.078 [+0.026, +0.134]}
\newcommand{\matchedFontanaSizeCallsAucVsX}{+0.023 [+0.006, +0.044]}
\newcommand{\matchedGateAucConfirm}{0.975}
\newcommand{\matchedGateAucExplore}{0.953}
\newcommand{\matchedGateAucVsXmConfirm}{+0.039 [+0.013, +0.063]}
\newcommand{\matchedGateAucVsXmExplore}{+0.059 [+0.024, +0.096]}
\newcommand{\matchedTokensAucConfirm}{0.963}
\newcommand{\matchedTokensAucVsXmConfirm}{+0.027 [$-$0.000, +0.054]}
\newcommand{\matchedTokensAucVsXmExplore}{+0.068 [+0.032, +0.104]}
\newcommand{\matchedCogAucConfirm}{0.936}
\newcommand{\matchedIntGateAucVsXmConfirm}{$-$0.081 [$-$0.145, $-$0.025]}
\newcommand{\matchedIntGateAucConfirm}{0.855}
\newcommand{\matchedWeightAucVsXmConfirm}{0.000 [$-$0.033, +0.029]}
\newcommand{\matchedGateVsXmConfirm}{+0.054 [$-$0.002, +0.105]}
\newcommand{\matchedGateVsXmExplore}{+0.064 [+0.011, +0.119]}
\newcommand{\matchedMethodRate}{14.1\%}
\newcommand{\matchedTokensVsXmConfirm}{+0.072 [+0.009, +0.133]}
\newcommand{\revMadeUtc}{20261026 235742 UTC}
\newcommand{\revShaShort}{321ac1632447}
\newcommand{\revLongReviews}{3{,}054}
\newcommand{\revBlobReviews}{3{,}518}
\newcommand{\revLongUnmapped}{303}
\newcommand{\revReviewers}{26}
\newcommand{\revNoCallLong}{3}
\newcommand{\revNoCallReviews}{300}
\newcommand{\revLeafTenRate}{0.4\%}
\newcommand{\revNonLeafTenRate}{1.0\%}
\newcommand{\revLeafTenN}{770}
\newcommand{\revNonLeafTenN}{823}
\newcommand{\revLeafTwentyFiveRate}{17.4\%}
\newcommand{\revNonLeafTwentyFiveRate}{24.3\%}
\newcommand{\revLeafTwentyFiveN}{230}
\newcommand{\revNonLeafTwentyFiveN}{571}
\newcommand{\revLeafFiftyRate}{80.0\%}
\newcommand{\revNonLeafFiftyRate}{81.6\%}
\newcommand{\revLeafFiftyN}{80}
\newcommand{\revNonLeafFiftyN}{365}
\newcommand{\revLeafHundredRate}{94.1\%}
\newcommand{\revNonLeafHundredRate}{94.0\%}
\newcommand{\revLeafHundredN}{17}
\newcommand{\revNonLeafHundredN}{151}
\newcommand{\revLeafOverRate}{85.7\%}
\newcommand{\revNonLeafOverRate}{95.0\%}
\newcommand{\revLeafOverN}{7}
\newcommand{\revNonLeafOverN}{40}
\newcommand{\revLongLeafMinusNonLeaf}{$-$0.033 [$-$0.141, +0.045]}
\newcommand{\revReviewersLeafLower}{4}
\newcommand{\revReviewersShown}{6}
\newcommand{\revDisagreeLongLeaf}{21.9\%}
\newcommand{\revDisagreeLongNonLeaf}{19.7\%}
\newcommand{\revDisagreeLongLeafN}{32}
\newcommand{\revDisagreeLongNonLeafN}{173}
\newcommand{\revCutoffReviewers}{10}
\newcommand{\revCutoffMin}{17}
\newcommand{\revCutoffMax}{34}
\newcommand{\revCutoffMedian}{21}
\newcommand{\revCutoffMccMin}{0.72}
\newcommand{\revCutoffMccMax}{0.91}
\newcommand{\revPosMedianMin}{30}
\newcommand{\revPosMedianMax}{44}
\newcommand{\revNegMedianMax}{15}
\newcommand{\revBlobHundredRate}{13.3\%}
\newcommand{\revBlobHundredN}{2{,}113}
\newcommand{\revBlobHundredDisagree}{38.0\%}
\newcommand{\revBlobTwoFiftyRate}{38.6\%}
\newcommand{\revBlobTwoFiftyN}{663}
\newcommand{\revBlobTwoFiftyDisagree}{61.7\%}
\newcommand{\revBlobFiveHundredRate}{48.9\%}
\newcommand{\revBlobFiveHundredN}{380}
\newcommand{\revBlobFiveHundredDisagree}{43.7\%}
\newcommand{\revBlobThousandRate}{43.1\%}
\newcommand{\revBlobThousandN}{246}
\newcommand{\revBlobThousandDisagree}{47.5\%}
\newcommand{\revBlobOverRate}{39.7\%}
\newcommand{\revBlobOverN}{116}
\newcommand{\revBlobOverDisagree}{41.7\%}
\newcommand{\revFlaggedLeafN}{100}
\newcommand{\revFlaggedRestN}{460}
\newcommand{\revFlaggedLeafMedianLen}{39}
\newcommand{\revFlaggedRestMedianLen}{42}
\newcommand{\revFlaggedLeafRate}{63.0\%}
\newcommand{\revFlaggedRestRate}{83.3\%}
\newcommand{\revBlobLargeCallsLowRate}{39.7\%}
\newcommand{\revBlobLargeCallsMidRate}{35.8\%}
\newcommand{\revBlobLargeCallsHighRate}{50.4\%}
\newcommand{\matchedGateAucVsTokensConfirm}{+0.011 [+0.002, +0.020]}
\newcommand{\matchedGatePctVsTokensConfirm}{+0.005 [$-$0.002, +0.013]}
\newcommand{\matchedGateMatchedVsTokensConfirm}{$-$0.018 [$-$0.056, +0.017]}
\newcommand{\matchedLinesPctVsTokensConfirm}{+0.005 [$-$0.002, +0.012]}
\newcommand{\matchedGatePctConfirm}{0.971}
\newcommand{\matchedLinesPctConfirm}{0.970}
\newcommand{\matchedTokensPctConfirm}{0.966}
\newcommand{\matchedCogPctConfirm}{0.934}
\newcommand{\matchedLinesOwnFlagsConfirm}{332}
\newcommand{\matchedTokensOwnFlagsConfirm}{252}
\newcommand{\matchedGateOwnFlagsConfirm}{157}
\newcommand{\matchedCogOwnFlagsConfirm}{110}
\newcommand{\matchedGateOwnMccConfirm}{0.744}
\newcommand{\matchedLinesOwnMccConfirm}{0.520}
\newcommand{\matchedGateAucVsTokensExplore}{$-$0.009 [$-$0.022, +0.003]}
\newcommand{\matchedGatePctVsTokensExplore}{$-$0.006 [$-$0.018, +0.003]}
\newcommand{\matchedGateMatchedVsTokensExplore}{$-$0.047 [$-$0.083, $-$0.010]}
\newcommand{\matchedLinesPctVsTokensExplore}{$-$0.002 [$-$0.011, +0.007]}
\newcommand{\matchedGatePctExplore}{0.954}
\newcommand{\matchedLinesPctExplore}{0.958}
\newcommand{\matchedTokensPctExplore}{0.960}
\newcommand{\matchedCogPctExplore}{0.892}
\newcommand{\matchedLinesOwnFlagsExplore}{358}
\newcommand{\matchedTokensOwnFlagsExplore}{265}
\newcommand{\matchedGateOwnFlagsExplore}{148}
\newcommand{\matchedCogOwnFlagsExplore}{124}
\newcommand{\matchedGateOwnMccExplore}{0.679}
\newcommand{\matchedLinesOwnMccExplore}{0.477}
\newcommand{\matchedOAucVsHvolConfirm}{$-$0.021 [$-$0.055, +0.014]}
\newcommand{\matchedSizeCallsAucVsHvolConfirm}{$-$0.008 [$-$0.032, +0.016]}
\newcommand{\logMadeUtc}{20261026 230649 UTC}
\newcommand{\logShaShort}{2fa130f412f6}
\newcommand{\logClassForms}{1813}
\newcommand{\logClassBetter}{64}
\newcommand{\logClassWorse}{308}
\newcommand{\logXFlagsRaw}{117}
\newcommand{\logXFlagsLog}{5}
\newcommand{\logXMccLog}{0.023}
\newcommand{\logOMccLog}{0.229}
\newcommand{\logOMinusRaw}{$-$0.043 [$-$0.094, $-$0.000]}
\newcommand{\logClassTopMcc}{0.295}
\newcommand{\logClassTopRaw}{0.273}
\newcommand{\logMethodForms}{280}
\newcommand{\logMethodBetter}{17}
\newcommand{\logMethodWorse}{100}
\newcommand{\logTokensFlagsRaw}{252}
\newcommand{\logTokensFlagsLog}{11}


\newcommand{\madz}{modified $z \geq 3.5$}
\newcommand{\fitwidth}[1]{\resizebox{\ifdim\width>\textwidth\textwidth\else\width\fi}{!}{#1}}
\newcommand{\sig}[1]{\texttt{#1}}
\newcommand{\Xs}{X$^*$}
\newcommand{\Xm}{X$^*_m$}
\newcommand{\Dflow}{D$_{\text{flow}}$}
\newcommand{\Oflow}{O$_{\text{flow}}$}
\newcommand{\Eflow}{E$_{\text{flow}}$}
\newcommand{\Aflow}{A$_{\text{flow}}$}
\newcommand{\dD}{$\Delta$D}
\newcommand{\dO}{$\Delta$O}
\newcommand{\dDO}{$\Delta$DO}

\title{\bfseries Does coupling divergence tell a god unit from a complex
algorithm?\\ A pre-registered test of weighting complexity concentration by
the data a unit hands on}

% Zenodo DOI, reserved before publishing and printed under the author block.
\newcommand{\doiline}{\ifdefined\thedoi\\\small\texttt{doi:\thedoi}\fi}

\author{Daniel Sagenschneider\\\small OfficeFloor\\\small\texttt{daniel@officefloor.net}\doiline}

\date{October 2026}

\begin{document}
\maketitle

\begin{abstract}
\noindent
A complexity outlier rule finds where complexity concentrates but, a
companion study found, cannot tell a god unit from a complex algorithm.
This paper pre-registers a test of one way to tell them apart: weight a
unit's complexity by its \emph{coupling divergence}, the data it hands on
minus the data it receives, so that a leaf is set aside. Coupling is
measured by resolving every call in \mlcqMeasured{} MLCQ repositories,
\seedThreeMeasured{} fresh Java repositories and \fontanaSystems{} Qualitas
Corpus systems, and the weighted signal is scored by the same scale-free
outlier rule as complexity alone against MLCQ's reviewers, Fontana's God
Class and five tool definitions. The primary, the pooled MCC of the
weighted form minus complexity alone against human Blob, is \primMicro{}:
no difference resolved, and underpowered. The exemption holds in two human
oracles of three, but the factor built on it loses to the exemption alone,
to a weighting by the number of methods, and to complexity alone on the
threshold-free ranking (AUC \primAucX{} against \primAucD{}). By the
pre-registered rule the construct's verdict is \emph{\consVerdict}: the
outgoing factor raises the MCC against Blob in the peak and density families
by \pOneOPeak{} and \pOneODensity{}. A post hoc threshold-free read
qualifies that: the outgoing factor flags two and a half times as many
classes and ranks Blobs no better than complexity alone (AUC
\matchedOAucConfirm{} against \matchedXAucConfirm{}), so its gain is the
flag rate. At method level length ranks Long Methods best of every form,
and MLCQ's individual reviews show why: the share calling a method a Long
Method rises from \revNonLeafTwentyFiveRate{} at 11 to 25 lines to
\revNonLeafFiftyRate{} at 26 to 50, each reviewer's best cutoff lies between
\revCutoffMin{} and \revCutoffMax{} lines, and at equal length a leaf is
called a Long Method no less often than a non-leaf. The reviewers' Long
Method is an absolute length, which a per-repository rule is the wrong
model for: a constant of \lwLengthCutoff{} lines or more, chosen on the
explore half, reaches MCC \lwLengthMccConfirm{} on the confirm half,
\lwLengthVsRule{} over the best per-repository form, and in a post hoc exploration
no form built from the method's variables, its data flow or the cost of
cutting it under that length does better.
\end{abstract}

\section{Introduction}

A companion study \cite{paperone} tested a god-class and god-method
detector with no raw-scale constant: a unit is flagged when its complexity
is a robust outlier for its own codebase, a modified z-score of 3.5 or more
on a scale fitted to that codebase \cite{iglewicz1993}. The rule found where
the complexity concentrates, and its extra flags, read in the source, were
small classes doing the work they were written to do. A complexity measure
cannot separate a god unit from a complex algorithm because the difference
is not in the complexity: both hold more of it than their neighbours. The
difference, that study concluded, is in the coupling, and it named this one
as the test of that reading.

The construct tested here is \emph{coupling divergence}. For a method with
$p$ parameters that makes calls into its own repository, each passing
$a(c)$ arguments, the data coming in is what the method receives and the
data going out is what it hands on; the divergence is the difference. A
method that receives more than it sends is a \emph{catchment}, a leaf that
processes what it is given, which is where the complexity of an algorithm
belongs. A method that sends more than it receives is a \emph{junction},
which creates and directs flow, and a unit that directs the data and also
holds the logic is what people call a god unit. The pre-registered claim is
that multiplying a unit's complexity by a factor that is zero for a
catchment and grows with the divergence of a junction recovers the units
people call god classes and methods better than complexity alone, by
setting complex leaves aside. The claim was committed, with its estimand,
decision rule, sample and bootstrap, before any coupling had been measured
on the confirmatory material.

The claim did not survive the test, though one part of it did. The primary
resolves no difference and is underpowered, and every resolved reading of
the whole factor goes against it. The exemption at the heart of the claim,
that a catchment is less often a god unit, is seen in the labels at method
level and in Fontana's God Class, and is unresolved for MLCQ's Blob; but the
factor built on it, which also amplifies a junction by its divergence, costs
more than the exemption gains. At the registered operating point the component that carries a signal is
the one the claim subtracted, the data a class sends out; a post hoc threshold-free read shows that gain to be a flag rate rather
than a better ranking. At method level the reviews themselves show what
people call a Long Method: a method of about twenty-five lines or more,
\lwLengthCutoff{} by the cutoff chosen on the explore half, leaf or not,
which no form built from calls, variables or data flow ranks better than
length does. This paper reports the
numbers in the order the plan fixed, the primary first, then the
diagnostic that tests the exemption on its own, then the consistency
analysis across oracles and definitions that reads the construct as a
whole, then the secondaries and sensitivities, and states where the
analysis as run departs from the plan as written.

\section{Background}

The relative outlier rule, the signal families and the oracles are the
companion study's and are summarised only as far as this paper needs them
\cite{paperone}. A unit is a method (every function the parser reports in a
\texttt{.java} file) or a class (the file). A method's lines are the code
lines of its own body; the methods of an anonymous class declared inside it
are units of their own, so a short method that returns a large anonymous
class counts as short. Ten method signals in three
families (size: lines, tokens, Halstead volume; structure: cyclomatic
\cite{mcCabe1976}, cognitive \cite{campbell2018}, nesting, exit points;
density: cyclomatic density, maintainability index, parameters) give 31
class signals by aggregating each as a sum, a mean and a maximum over the
file's methods, plus the number of methods. The rule standardises a signal
within each repository by its median and median absolute deviation and
flags units at \madz{}; where more than half the units tie at the median
the MAD is zero and the rule falls back to the mean absolute deviation, a
detail that matters here because the weighted forms are zero on every
exempt unit.

Coupling as an ingredient of a god definition is old. Lanza and Marinescu's
God Class conjoins a weighted method count, low cohesion and access to
foreign data \cite{lanza2006, marinescu2004}, and PMD implements it, so
agreement with PMD is partly circular for any coupling measure. Henry and
Kafura's information flow multiplies a procedure's length by the square of
fan-in times fan-out \cite{henry1981}, and Kitchenham, Pickard and Linkman's
replication found size predicted as well \cite{kitchenham1990}, which is
why a size-only weighting is one of this paper's arms. Divergence differs
from information flow in subtracting what comes in from what goes out, so
that a unit high on both is not scored highest; whether that subtraction
adds anything is the question the consistency analysis answers. Yourdon and
Constantine's tramp data, data passed through modules that do not use it
\cite{yourdon1979}, is the ancestor of the construct behind divergence, a
\emph{sinew}: data that reaches the code that needs it only by riding the
call chain. Tempero and Ralph ask that a metric name its construct, entity
and attribute and show convergent and discriminant validity before it is
used for a decision \cite{tempero2026metrics}; in those terms the construct
is the sinew, the entity a method or class, the attribute the direction of
the data coupled to the control flow, and divergence one indicator of it.
Arif, Kuutila and Ralph's factor analysis of class metrics found in-coupling
and out-coupling loading as separate constructs \cite{arif2026}, which is
the reason the outgoing side is measured and read on its own here, as a
pre-registered secondary. Zhang, Wen and Tempero's typed dependency
distributions, on tool-labelled smells, found data access the dependency
kind that most separates smelly from clean code and labelled the God Class a
strong method consumer \cite{zhang2025deps}; theirs is the closest published
typology to the in and out counted here, with direction taken relative to a
pair of smells rather than to the unit's own flow.

Two older measures sit beside the forms the exploration in
Section~\ref{sec:explore-forms} arrives at. Chidamber and Kemerer's
response for a class is the number of methods a class declares plus the
number it calls \cite{ck1994}, size and out-calls in one sum; and Lanza and
Marinescu's method-level Intensive and Dispersed Coupling flag a method by
the number of methods it calls, conjoined with nesting
\cite{lanza2006}. Their Brain Method is long, branchy, deeply nested and
touches many variables, with no call count in it. El Emam and colleagues
showed that object-oriented metrics lose most of their predictive power once
class size is controlled \cite{elemam2001}, which is why every form here is
read against a size-only competitor. Cognitive complexity counts a whole
\texttt{switch} once so that a lookup table is not scored as branchy
\cite{campbell2018}, the one place in the detection literature where a long
self-contained method is deliberately discounted; a search found no
detection rule that excludes methods which make no call.

The human oracle is MLCQ \cite{mlcq2020}, professional developers' reviews
of Java open-source classes and methods, with the companion study's
strict-majority aggregation and its explore and confirm halves frozen before
any repository was measured. Classifiers trained on MLCQ reach a median MCC
of about 0.51 for Blob and 0.81 for Long Method \cite{mlcq2023}; the rules
here see no label, so those are a reference point, not a baseline. Fontana,
M\"antyl\"a, Zanoni and Marino's God Class instances on the Qualitas Corpus
\cite{fontana2016, tempero2010} are a second human oracle from a different
team on older code, read only as a paired difference because their
instances were found by tools before they were checked by hand and are
balanced and extreme, so any detector separates them more easily than in
real code \cite{dinucci2018}.

\section{The claim and how it is read}
\label{sec:claim}

\paragraph{The measures.} For each method, under the primary \emph{flow}
definition of coupling, $\mathit{in} = 1 + p$ and $\mathit{out} = \sum_c
(1 + a(c))$ over the method's internal calls: the call itself is control
flow, and each parameter or argument is data coupled to it. Two secondary
definitions run through every analysis: \emph{data}, parameters and
arguments only, and \emph{message}, which treats the first item as the
message itself. An internal call is one the symbol solver resolves to a
declaration in the repository's own sources; a call to the JDK or to a
dependency is external and counts for nothing. The divergence is $d =
\mathit{out} - \mathit{in}$. A class's $\mathit{out}$ sums its methods'
calls that leave the file, and its $\mathit{in}$ sums the parameters of its
methods callable from outside the file, non-private methods in the primary
variant and public methods in a secondary; constructors are structure, not
inputs, and are counted in a reported variant only.

\paragraph{The forms.} The divergence factor is $f_D = \max(d + 1, 0)$:
a catchment ($d < 0$) is 0, a pipe ($d = 0$) is 1, a junction with
divergence $j$ is $j + 1$. It is the product of two parts that are
measured apart, the exemption $f_E = [d \geq 0]$ and the amplification
$f_A = \max(d + 1, 1)$. The outgoing factor is $f_O = \mathit{out} + 1$,
registered as a secondary after an exploratory observation on agent-evolved
code that the outgoing side carried the signal. The size null is $f_N$,
the class's number of methods, which carries no coupling. Every class signal
X gives the forms X$\cdot$D, X$\cdot$E, X$\cdot$A, X$\cdot$O and
X$\cdot$N; every one is scored by the same rule on every unit of the
repository, exempt units as zeros, exactly as complexity alone.

\paragraph{The primary.} The base signal \Xs{} is chosen on the MLCQ explore
half by complexity alone, among the peak and density families (never size,
never parameters), as the signal whose flags maximise the pooled MCC
against human Blob; the multiplied form is not consulted. The estimand is
the pooled MCC of \Xs$\cdot$\Dflow{}, class form A with non-private inputs,
minus \Xs{}, both by \madz{}, against MLCQ human Blob over the reviewed
classes of the confirm repositories, as a paired difference with a 95\%
percentile bootstrap over repositories (\bootB{} resamples, seed
\bootSeed{}). A lower bound above 0 reads \emph{the divergence factor
refines complexity}; an upper bound below 0, \emph{it makes it worse};
otherwise \emph{no difference resolved}. An effective-sample gate, sized on
the explore half before confirm was read, marks the verdict
\emph{underpowered} when fewer confirm repositories carry a reviewed Blob
than the gate requires.

\paragraph{The mechanism and the size null.} If the primary reads
\emph{refines}, its mechanism is read from the exemption-only arm and the
leaf diagnostic: the refinement is by the exemption, which the claim names,
only if \Xs$\cdot$\Eflow{} also refines and catchments are not called Blob
more often than the rest. Beside it, \Xs$\cdot$\Dflow{} is read against
\Xs$\cdot$N: a lower bound above 0 says the coupling factor is
\emph{beyond size}; an upper bound below 0, that the size factor beats it.

\paragraph{The leaf diagnostic.} Among the units \Xs{} flags, the human
Blob rate of catchments against the rest, for each definition and input
variant. The exemption is \emph{supported} if catchments are called Blob
less often (upper bound below 0), \emph{contradicted} if more often,
otherwise \emph{unresolved}. It is reported whatever the primary says.

\paragraph{The construct.} The consistency analysis asks whether divergence
holds up across ways of defining a god unit and ways of calculating
complexity. For every oracle, level, signal family, relation (whether the
signal is an ingredient of the oracle's definition) and coupling definition,
three paired differences are averaged over the family's signals: \dD{}
$= \mathrm{MCC}(\mathrm{X}\cdot\mathrm{D}) - \mathrm{MCC}(\mathrm{X})$,
\dO{} $= \mathrm{MCC}(\mathrm{X}\cdot\mathrm{O}) - \mathrm{MCC}(\mathrm{X})$
and \dDO{} $= \mathrm{MCC}(\mathrm{X}\cdot\mathrm{D}) -
\mathrm{MCC}(\mathrm{X}\cdot\mathrm{O})$, the last asking whether what comes
in matters. One bootstrap resamples repositories once per replicate and
recomputes every cell, so the cells share their intervals' dependence and
are never counted as independent evidence. Ingredient cells and PMD
GodClass cells (which contain coupling) are reported and not counted. Four
predictions were registered on the flow definition: P1, \dD{} $> 0$ for
human Blob in the peak and density families; P2, \dD{} larger in peak and
structure families than in size; P3, \dDO{} $\geq 0$ for human Blob in the
counted families, against the competing expectation \dDO{} $\leq 0$ from the
exploratory observation; P4, \dD{} positive in a majority of the counted
tool cells. The verdict on the construct is \emph{supported} if P1 holds
and P3 does not read outgoing-only; \emph{outgoing data only} if P1 holds
for \dO{} and P3 reads outgoing-only; \emph{not supported} if neither \dD{}
nor \dO{} has a lower bound above 0 for human Blob in any non-size family;
otherwise \emph{unresolved}.

\paragraph{What was seen before the plan.} The companion study's main
confirmatory result against PMD, which motivates the question, had been
read; its human-label read-out had not. Part of its MLCQ step output had
been seen while checking that the run was progressing, none of it a result
against the human labels. The exploratory observation that motivated the
outgoing-only arm came from the author's agent-evolved petclinic runs. The
plan was amended ten times after registration, each amendment a commit
stating what had been seen. The first nine were made before any confirm or
seed-3 repository had been measured. The tenth corrects a fault in the
harness's populations found after the confirmatory results had been read;
it is stated in Section~\ref{sec:deviations} and in full in the ancillary
data.

\section{Study design}
\label{sec:design}

\paragraph{Material.} Three measured sets. (1) The \mlcqRepos{} MLCQ
repositories at their labelled commits, \mlcqMeasured{} of them measured,
\mlcqMovedMeasured{} at a new GitHub location verified to hold the labelled
commit; the confirm half holds \mlcqConfirmMeasured{} repositories, of
which \mlcqConfirmUsable{} measured and pass the companion study's
substance floor. MLCQ holds \mlcqBlobReviewed{} reviewed classes with
\mlcqBlob{} Blob positives and \mlcqLongReviewed{} reviewed methods with
\mlcqLong{} Long Method positives after strict-majority aggregation, over
both halves. (2) A fresh sample of \seedThreeMeasured{} usable Java
repositories drawn with seed 3 by the companion study's sampler, stratified
by popularity tier, after removing every repository that appears in the
companion study's pilot, its seed-2 sample or MLCQ (\seedThreeExcluded{}
of the \seedThreeDrawUsable{} usable in the draw); \seedThreeUsable{} pass
the substance floor. It carries the tool oracles: PMD GodClass, an
approximate DECOR Blob composed from CK's counts \cite{moha2010, cktool},
WMC $\geq 47$, PMD NPath at method level, and an internal Brain Method.
(3) Fontana's God Class: \fontanaInstances{} manually evaluated instances,
\fontanaGodInstances{} of them God Classes, on \fontanaSystems{} system
versions of the Qualitas Corpus \cite{tempero2010}, measured from the
corpus's source archives rather than from git; several instances share a
file, so they land on \fontanaFilesReviewed{} reviewed files with
\fontanaFilesGod{} positives, of which \fontanaFilesReviewedMeasured{} and
\fontanaFilesGodMeasured{} land on a measured file. The corpus's maintainer
agreed to its use. Test directories, build output and generated sources
are dropped before any rule is fitted, matched against each file's source
root so that a package named \texttt{build} or \texttt{target} is kept.

\paragraph{Two result sets.} Complexity is the companion study's extractor
at its released commit, so every baseline number is the number that study
would produce on the same unit; coupling is a JavaParser measurer with its
symbol solver, resolving against the repository's sources and the JDK with
no build and no dependencies, under a time budget of eight hours per
repository. The two are joined by file and start line (the \emph{joined}
set, which decides). As a robustness set, every complexity field was
re-implemented on JavaParser's syntax tree so that complexity and coupling
come from one parser and no method is lost to the complexity extractor's
blind spots (the \emph{JavaParser-only} set), and every verdict is marked
agrees or differs.

\paragraph{The matrix.} Every form of every signal is scored by every rule
of the companion study at its default constant; the paper reads \madz{} and
reports the rest in the ancillary analysis file. Against human Blob the
class matrix holds \matrixBlobRows{} form-and-rule rows; against Long
Method, \matrixLongRows{}. Beside the paired differences, an exploratory
intersection flags the outliers of complexity and of divergence separately
by the same rule and reports the positives by region.

\section{Results}

\subsection{Diagnostics, read first}
\label{sec:diag}

\paragraph{The join and the measurer.} On the MLCQ confirm half the
measurer's methods joined \joinMlcqShare{} of the \joinMlcqLizardUnits{}
units the complexity extractor reports, and found \joinMlcqMeasurerExtra{}
methods (\joinMlcqMeasurerExtraShare{} of its own) the extractor does not,
the expected loss after some class-level field initialisers. Of the
\joinMlcqClassUnits{} class units \joinMlcqClassCoupledShare{} carry
coupling, among them every one of the \uncoupledBlobReviewed{} reviewed Blob
classes; \uncoupledLongUncoupled{} of the \uncoupledLongReviewed{} reviewed
Long Method units carry none, none of them positive; on Fontana
\uncoupledFontanaUncoupled{} reviewed classes carry none,
\uncoupledFontanaPositiveUncoupled{} of them positives.
On the explore half the JavaParser re-implementation of the complexity
fields agrees with the extractor on \agreeUnits{} joined methods at exact
matches of \agreeCcExact{} for cyclomatic complexity, \agreeNlocExact{} for
lines, \agreeParamsExact{} for parameters, \agreeCognitiveExact{} for
cognitive complexity (pooled Spearman \agreeCognitiveRho{}) and
\agreeNestingExact{} for nesting (\agreeNestingRho{}); Halstead volume and
the maintainability index match exactly on \agreeHvolExact{} and
\agreeMiExact{} of methods, floating-point differences, at Spearman
\agreeHvolRho{} and \agreeMiRho{}. The measurer's cognitive complexity
agrees with PMD's at a median per-repository Spearman of \pmdCogMeasurer{}
against the extractor's \pmdCogPaperOne{}.

\paragraph{Resolution.} Over the confirm half \callsMlcqInternal{} of
calls resolved to the repository's own code, \callsMlcqJdk{} to the JDK,
\callsMlcqUnresolved{} not at all (almost always a dependency's type), and
\callsMlcqShouldHave{} to a receiver declared in the repository whose
method could not be resolved; no repository spent its time budget, on any
set. On the seed-3 sample the shares are \callsSeedThreeInternal{},
\callsSeedThreeJdk{}, \callsSeedThreeUnresolved{} and
\callsSeedThreeShouldHave{}; on Fontana's older code, with fewer
dependencies, \callsFontanaInternal{} internal and
\callsFontanaUnresolved{} unresolved. A call that is not resolved is
treated as external and adds nothing to $\mathit{out}$, so $\mathit{out}$
reads low in proportion. \diMlcqRepos{} of the \mlcqConfirmUsable{} usable
confirm repositories declare a dependency-injection framework, whose
injected collaborators are calls the parser cannot see; the primary is
re-read without them in Section~\ref{sec:sens}. Of the data a class sends
out under flow, \throwNewMlcqClassShare{} is the construction of an
exception to throw, counted as registered.

\paragraph{Roles.} Under flow on the confirm half, \roleMlcqClassCatchmentUnits{}
of classes are catchments holding \roleMlcqClassCatchmentCc{} of the
cyclomatic complexity, \roleMlcqClassPipeUnits{} are pipes and
\roleMlcqClassJunctionUnits{} are junctions holding
\roleMlcqClassJunctionCc{}. At method level \roleMlcqMethodCatchmentUnits{}
of methods are catchments, every method with no internal call among them,
holding \roleMlcqMethodCatchmentCc{} of the complexity, and
\roleMlcqMethodJunctionUnits{} are junctions holding
\roleMlcqMethodJunctionCc{}. Under the data definition, which does not see a
call with no arguments, \roleMlcqMethodDataPipeUnits{} of methods are pipes
rather than catchments, as the plan anticipated of a method with no
parameters.

\paragraph{Zeros and the scale.} Exempt units are zeros of the weighted
form, and they are more than half the units in the median repository:
\zeroMlcqDClass{} of classes under \Xs$\cdot$\Dflow{}, so the rule fell to
the mean-absolute-deviation fallback in \scaleMlcqDMeanAd{} of the
\scaleMlcqClassRepos{} confirm repositories, against \scaleMlcqXMad{} on the
MAD scale for \Xs{} alone. At method level \Xm{} itself is zero-heavy
(\zeroMlcqXm{} zeros in the median repository, the fallback in
\scaleMlcqXmMeanAd{} of \scaleMlcqMethodRepos{}), and the weighted form has
\zeroMlcqDmMethod{} zeros and uses the fallback everywhere. The two arms of
the primary are therefore read under different scale estimators, which is
why the arm fitted on the non-exempt units only is reported as a sensitivity.

\subsection{The explore-half choices}
\label{sec:explore}

On the \xClassRepos{} explore repositories carrying a reviewed class,
complexity alone chose \sig{\xClass} as \Xs{} at a pooled MCC of
\xClassMcc{}, ahead of \sig{\xClassSecond} at \xClassSecondMcc{} and
\sig{\xClassThird} at \xClassThirdMcc{}; the companion study's
\sig{cognitive\_max}, chosen there against PMD, ranks \xClassCogMaxRank{}th
of \xClassCandidates{} at \xClassCogMaxMcc{}. At method level, on
\xMethodRepos{} repositories, \sig{\xMethod} was chosen at \xMethodMcc{}
ahead of \sig{\xMethodSecond} at \xMethodSecondMcc{}. The explore half's own
primary estimate was \gateExploreEstimate{}, a half-width of
\gateHalfWidth{} over \gateExplorePositiveRepos{} repositories, so the gate
set \gateRequired{} confirm repositories with a reviewed Blob for a
half-width of 0.10 (the gate as first registered, from the per-repository
standard deviation \gateVOneS{}, gives \gateVOneRequired{} and does not
decide). The confirm half carries \primGateReposWithBlob{}: the primary is
underpowered by the plan's own rule, whatever it reads.

\subsection{The primary}
\label{sec:primary}

\begin{table}[t]
\centering
\caption{The primary read-out against MLCQ human Blob on the confirm half:
\primUniv{} reviewed classes, \primPos{} of them Blob, in \primRepos{}
repositories. Every arm is scored by \madz{} on every class of its
repository. Flagged and hits are pooled counts over the reviewed classes;
the paired difference is the pooled MCC of the arm minus that of \Xs{}, with
its 95\% percentile bootstrap interval over repositories.}
\label{tab:primary}
\small
\setlength{\tabcolsep}{4pt}
\fitwidth{% generated by paper/scripts/tables.py; do not edit
\begin{tabular}{lrrrrrrl}
\toprule
Arm (class form A, flow, non-private inputs) & MCC & Precision & Recall & FPR & Flagged & Hits & Paired difference against X$^*$ (95\% CI) \\
\midrule
X$^*$ = \texttt{nloc\_max}, complexity alone & 0.187 & 0.26 & 0.31 & 0.104 & 117 & 30 &  \\
X$^*\cdot$D$_{\text{flow}}$, the primary arm & 0.146 & 0.19 & 0.41 & 0.208 & 214 & 40 & $-$0.041 [$-$0.130, +0.054] \\
X$^*\cdot$E$_{\text{flow}}$, exemption only & 0.212 & 0.26 & 0.37 & 0.122 & 138 & 36 & +0.025 [$-$0.052, +0.100] \\
X$^*\cdot$A$_{\text{flow}}$, amplification only & 0.196 & 0.20 & 0.56 & 0.268 & 279 & 55 & +0.009 [$-$0.075, +0.104] \\
X$^*\cdot$N, the size null & 0.236 & 0.22 & 0.60 & 0.252 & 270 & 59 & +0.049 [$-$0.023, +0.121] \\
X$^*\cdot$D$_{\text{flow}}$ with constructors as inputs & 0.125 & 0.20 & 0.27 & 0.124 & 130 & 26 & $-$0.062 [$-$0.130, +0.011] \\
X$^*\cdot$O$_{\text{flow}}$, outgoing factor only & 0.272 & 0.23 & 0.69 & 0.279 & 301 & 68 & +0.085 [+0.001, +0.182] \\
\bottomrule
\end{tabular}
}
\end{table}

Table~\ref{tab:primary} gives every arm. Complexity alone,
\sig{\xClass} at \madz{}, flags \primXFlagged{} of the \primUniv{} reviewed
confirm classes and hits \primXTp{} of the \primPos{} Blobs: MCC
\primXMcc{}, precision \primXPrecision{}, recall \primXRecall{}. The primary
arm flags \primDFlagged{} and hits \primDTp{}: MCC \primDMcc{}, precision
\primDPrecision{}, recall \primDRecall{}. The paired difference is
\textbf{\primMicro{}}, and with \primGateReposWithBlob{} repositories
against a gate of \gateRequired{} the verdict is \textbf{\primVerdict}. The
confirm half-width is \primHalfWidthConfirm{}, as the explore half
predicted. The macro difference, one repository one vote, is
\primMacro{}. The point estimate is on the wrong side, and the
hypothesised shape, higher precision at a small cost in recall, is
inverted: the factor gains recall (\primDRecall{} against \primXRecall{})
and loses precision (\primDPrecision{} against \primXPrecision{}), the
false positives rising from \primXFp{} to \primDFp{}. The threshold-free
comparison resolves what the operating point does not: pooled over the
reviewed classes, the modified z of \Xs{} ranks Blobs above non-Blobs at an
AUC of \primAucX{} and the weighted form at \primAucD{}, a paired difference
of \primAucDiff{}, complexity alone the better ranking.

\paragraph{With constructors as inputs.} Counting non-private constructors
in a class's inputs gives \primCtorMicro{}, \primCtorReading{}, which
\primCtorAgrees{} with the primary; the arm's MCC is \primCtorMcc{} with
precision \primCtorPrecision{} and recall \primCtorRecall{}, and the
difference between the two arms is \primMethodsVsCtorMicro{}.

\paragraph{The mechanism is not read.} The plan reads the mechanism of a
refinement only when the primary reads \emph{refines}, so it is not read.
The two halves of the factor are nonetheless informative. Exemption only,
\Xs$\cdot$\Eflow{}, is \primEMicro{} against \Xs{} (MCC \primEYMcc{},
precision \primEYPrecision{}, recall \primEYRecall{}); amplification only,
\Xs$\cdot$\Aflow{}, is \primAmpMicro{} (MCC \primAmpYMcc{}, recall
\primAmpYRecall{} at precision \primAmpYPrecision{}). Neither is resolved
against complexity alone, but the whole factor is resolved against the
exemption alone: \primDvsEMicro{}, \primDvsEReading{}. Against the
amplification alone it is \primDvsAMicro{}, \primDvsAReading{} on the joined
set and \primJpDvsAMicro{} on the JavaParser-only set, where it reads
\primJpDvsAReading{}. Multiplying the two halves together does worse than
either half on its own: resolved against the exemption on both sets, and
against the amplification on one.

\paragraph{The size null.} Weighting \Xs{} by the number of methods, which
carries no coupling, gives \primNvsXMicro{} against \Xs{} (MCC
\primNMcc{}, precision \primNPrecision{}, recall \primNRecall{}), not
resolved on the pooled MCC and \primNvsXMacro{} on the macro. Against it the
coupling factor is \textbf{\primDvsNMicro{}}: \primDvsNReading{}. The
rival explanation the plan recorded for any class-level gain, that a Blob is
a big class with one complex method, is not merely left standing: the size
weighting beats the coupling weighting by a resolved margin.

\paragraph{The JavaParser-only set.} On \primJpUniv{} reviewed classes the
primary is \primJpMicro{}, \primJpVerdict{}; the AUCs are \primJpAucX{}
and \primJpAucD{} (\primJpAucDiff{}); the halves read \primJpEMicro{} and
\primJpAmpMicro{}, the whole factor against them \primJpDvsEMicro{} and
\primJpDvsAMicro{}, and against the size null \primJpDvsNMicro{}. Every
verdict agrees.

\subsection{The leaf diagnostic}
\label{sec:leaf}

\begin{table}[t]
\centering
\caption{The leaf diagnostic: among the reviewed units \Xs{} (or \Xm{})
flags, how many are catchments ($d < 0$), how many of those the reviewers
call positive, the same for the rest, and the paired difference in positive
rate (catchments minus the rest) with its interval over repositories. The
exemption is supported by an upper bound below 0 and contradicted by a
lower bound above 0.}
\label{tab:leaf}
\small
\setlength{\tabcolsep}{4pt}
\fitwidth{% generated by paper/scripts/tables.py; do not edit
\begin{tabular}{llrrrrll}
\toprule
Oracle & Definition, inputs & Catchments flagged & of them positive & Rest flagged & of them positive & Difference in positive rate (95\% CI) & Reading \\
\midrule
MLCQ Blob & flow, non-private & 38 & 8 & 79 & 22 & $-$0.068 [$-$0.238, +0.102] & unresolved \\
MLCQ Blob & flow, public & 33 & 7 & 84 & 23 & $-$0.062 [$-$0.229, +0.118] & unresolved \\
MLCQ Blob & flow, non-private, with constructors & 41 & 8 & 76 & 22 & $-$0.094 [$-$0.260, +0.069] & unresolved \\
MLCQ Blob & flow, public, with constructors & 34 & 7 & 83 & 23 & $-$0.071 [$-$0.236, +0.105] & unresolved \\
MLCQ Blob & data, non-private & 44 & 8 & 73 & 22 & $-$0.120 [$-$0.284, +0.041] & unresolved \\
MLCQ Blob & data, public & 34 & 8 & 83 & 22 & $-$0.030 [$-$0.196, +0.145] & unresolved \\
MLCQ Blob & data, non-private, with constructors & 52 & 10 & 65 & 20 & $-$0.115 [$-$0.280, +0.047] & unresolved \\
MLCQ Blob & data, public, with constructors & 44 & 10 & 73 & 20 & $-$0.047 [$-$0.213, +0.122] & unresolved \\
MLCQ Blob & message, non-private & 37 & 7 & 80 & 23 & $-$0.098 [$-$0.266, +0.060] & unresolved \\
MLCQ Blob & message, public & 31 & 5 & 86 & 25 & $-$0.129 [$-$0.296, +0.031] & unresolved \\
MLCQ Blob & message, non-private, with constructors & 38 & 7 & 79 & 23 & $-$0.107 [$-$0.278, +0.049] & unresolved \\
MLCQ Blob & message, public, with constructors & 33 & 6 & 84 & 24 & $-$0.104 [$-$0.269, +0.054] & unresolved \\
MLCQ Long Method & flow, all calls & 28 & 13 & 82 & 63 & $-$0.304 [$-$0.541, $-$0.082] & supported (catchments are called positive less often) \\
MLCQ Long Method & data, all calls & 28 & 15 & 82 & 61 & $-$0.208 [$-$0.413, $-$0.017] & supported (catchments are called positive less often) \\
MLCQ Long Method & message, all calls & 27 & 12 & 83 & 64 & $-$0.327 [$-$0.557, $-$0.102] & supported (catchments are called positive less often) \\
Fontana God Class & flow, non-private & 20 & 15 & 98 & 93 & $-$0.199 [$-$0.362, $-$0.001] & supported (catchments are called positive less often) \\
Fontana God Class & data, non-private & 19 & 15 & 99 & 93 & $-$0.150 [$-$0.316, +0.050] & unresolved \\
Fontana God Class & message, non-private & 22 & 17 & 96 & 91 & $-$0.175 [$-$0.322, 0.000] & unresolved \\
\bottomrule
\end{tabular}
}
\end{table}

The exemption is the claim in the question, and this is its own test,
reported whatever the primary says (Table~\ref{tab:leaf}). Against MLCQ Blob
it is \leafReading{}: of the \primXFlagged{} reviewed classes \Xs{} flags on
confirm, under flow with non-private inputs, \leafCatchFlagged{} are
catchments and \leafCatchPositive{} of them are Blobs, \leafCatchRate{}; of
the \leafRestFlagged{} others, \leafRestPositive{} are Blobs,
\leafRestRate{}; the difference is \leafDiff{}, and it is unresolved in
\leafClassUnresolved{} of the \leafClassVariants{} definition and input
variants, every point estimate on the exemption's side. Against MLCQ Long
Method it is \leafMethodReading{}: \leafMethodCatchPositive{} of the
\leafMethodCatchFlagged{} flagged catchments are Long Methods,
\leafMethodCatchRate{} against \leafMethodRestRate{} for the rest,
\leafMethodDiff{}, and the same under data (\leafMethodDataDiff{}) and
message (\leafMethodMessageDiff{}). Against Fontana's God Class it is
\leafFontanaReading{} under flow: \leafFontanaCatchRate{} of the flagged
catchments are God Classes against \leafFontanaRestRate{} of the rest,
\leafFontanaDiff{}, supported in \leafFontanaFlowSupported{} of 4 flow variants and
\leafFontanaSupported{} of \leafFontanaVariants{} variants in all.
A complex unit that receives more than it sends is called a god unit less
often by two of the three human oracles, and not detectably less often by
the third. The exemption is a distinction the labels carry among the units complexity
flags; Section~\ref{sec:explore-forms} reads it against the individual
reviews, where it holds for the branchy leaves complexity flags and not for
leaves at large, and the rest of the results say why the factor that rests
on it does not improve detection.

\subsection{The construct: the consistency analysis}
\label{sec:cons}

\begin{table}[t]
\centering
\caption{The consistency cells of the human oracles under the flow
definition, class form A with non-private inputs, every signal of the family
averaged: \dD{} is the paired MCC difference of X$\cdot$D against X, \dO{}
of X$\cdot$O against X, and \dDO{} of X$\cdot$D against X$\cdot$O. One
bootstrap over repositories gives every interval. $n$ is the number of
repositories (systems) carrying the oracle.}
\label{tab:cells-human}
\small
\setlength{\tabcolsep}{4pt}
\fitwidth{% generated by paper/scripts/tables.py; do not edit
\begin{tabular}{llllrlll}
\toprule
Oracle & Family & Relation & Counted & $n$ & $\Delta$D (95\% CI) & $\Delta$O (95\% CI) & $\Delta$DO (95\% CI) \\
\midrule
MLCQ Blob (class) & size & non-ingredient & yes & 208 & $-$0.032 [$-$0.114, +0.048] & +0.048 [$-$0.014, +0.108] & $-$0.080 [$-$0.144, $-$0.014] \\
MLCQ Blob (class) & peak & non-ingredient & yes & 208 & +0.041 [$-$0.031, +0.114] & +0.125 [+0.054, +0.195] & $-$0.084 [$-$0.139, $-$0.029] \\
MLCQ Blob (class) & density & non-ingredient & yes & 208 & +0.094 [+0.013, +0.174] & +0.151 [+0.076, +0.227] & $-$0.058 [$-$0.110, $-$0.009] \\
MLCQ Long Method (method) & size & non-ingredient & yes & 202 & $-$0.232 [$-$0.317, $-$0.144] & $-$0.159 [$-$0.200, $-$0.121] & $-$0.073 [$-$0.156, +0.012] \\
MLCQ Long Method (method) & structure & non-ingredient & yes & 202 & $-$0.044 [$-$0.115, +0.028] & $-$0.028 [$-$0.086, +0.034] & $-$0.017 [$-$0.061, +0.024] \\
MLCQ Long Method (method) & density & non-ingredient & yes & 202 & +0.280 [+0.197, +0.362] & +0.222 [+0.162, +0.280] & +0.058 [+0.006, +0.110] \\
Fontana God Class (class) & size & non-ingredient & yes & 72 & $-$0.073 [$-$0.193, +0.037] & +0.063 [$-$0.004, +0.121] & $-$0.136 [$-$0.243, $-$0.038] \\
Fontana God Class (class) & peak & non-ingredient & yes & 72 & +0.169 [+0.068, +0.269] & +0.315 [+0.249, +0.381] & $-$0.146 [$-$0.230, $-$0.067] \\
Fontana God Class (class) & density & non-ingredient & yes & 72 & +0.413 [+0.326, +0.503] & +0.573 [+0.499, +0.658] & $-$0.160 [$-$0.247, $-$0.079] \\
\bottomrule
\end{tabular}
}
\end{table}

\begin{table}[t]
\centering
\caption{The same cells for the tool oracles on the seed-3 sample and for
the internal definitions. Ingredient cells and PMD GodClass cells are
reported and not counted toward the construct.}
\label{tab:cells-tool}
\small
\setlength{\tabcolsep}{4pt}
\fitwidth{% generated by paper/scripts/tables.py; do not edit
\begin{tabular}{llllrlll}
\toprule
Oracle & Family & Relation & Counted & $n$ & $\Delta$D (95\% CI) & $\Delta$O (95\% CI) & $\Delta$DO (95\% CI) \\
\midrule
PMD GodClass (class) & size & ingredient & no & 490 & $-$0.285 [$-$0.321, $-$0.244] & $-$0.125 [$-$0.150, $-$0.099] & $-$0.159 [$-$0.193, $-$0.123] \\
PMD GodClass (class) & size & non-ingredient & no & 490 & $-$0.273 [$-$0.309, $-$0.230] & $-$0.142 [$-$0.160, $-$0.124] & $-$0.131 [$-$0.163, $-$0.096] \\
PMD GodClass (class) & peak & ingredient & no & 490 & +0.025 [$-$0.003, +0.056] & +0.148 [+0.123, +0.172] & $-$0.123 [$-$0.153, $-$0.091] \\
PMD GodClass (class) & peak & non-ingredient & no & 490 & $-$0.043 [$-$0.073, $-$0.009] & +0.060 [+0.041, +0.081] & $-$0.103 [$-$0.133, $-$0.072] \\
PMD GodClass (class) & density & ingredient & no & 490 & +0.190 [+0.161, +0.221] & +0.295 [+0.268, +0.320] & $-$0.105 [$-$0.135, $-$0.072] \\
PMD GodClass (class) & density & non-ingredient & no & 490 & +0.156 [+0.129, +0.187] & +0.241 [+0.216, +0.264] & $-$0.085 [$-$0.112, $-$0.056] \\
DECOR Blob (class) & size & ingredient & no & 490 & $-$0.302 [$-$0.348, $-$0.259] & $-$0.108 [$-$0.144, $-$0.075] & $-$0.194 [$-$0.223, $-$0.163] \\
DECOR Blob (class) & size & non-ingredient & yes & 490 & $-$0.241 [$-$0.281, $-$0.199] & $-$0.085 [$-$0.104, $-$0.066] & $-$0.157 [$-$0.187, $-$0.121] \\
DECOR Blob (class) & peak & non-ingredient & yes & 490 & $-$0.017 [$-$0.049, +0.017] & +0.098 [+0.082, +0.116] & $-$0.115 [$-$0.144, $-$0.082] \\
DECOR Blob (class) & density & non-ingredient & yes & 490 & +0.129 [+0.085, +0.172] & +0.217 [+0.189, +0.246] & $-$0.088 [$-$0.117, $-$0.057] \\
WMC $\geq 47$ (class) & size & ingredient & no & 491 & $-$0.356 [$-$0.409, $-$0.300] & $-$0.158 [$-$0.191, $-$0.124] & $-$0.199 [$-$0.236, $-$0.159] \\
WMC $\geq 47$ (class) & size & non-ingredient & yes & 491 & $-$0.337 [$-$0.382, $-$0.278] & $-$0.166 [$-$0.192, $-$0.140] & $-$0.170 [$-$0.206, $-$0.127] \\
WMC $\geq 47$ (class) & peak & ingredient & no & 491 & +0.038 [0.000, +0.077] & +0.178 [+0.151, +0.204] & $-$0.140 [$-$0.173, $-$0.104] \\
WMC $\geq 47$ (class) & peak & non-ingredient & yes & 491 & $-$0.039 [$-$0.073, +0.001] & +0.088 [+0.065, +0.111] & $-$0.127 [$-$0.159, $-$0.089] \\
WMC $\geq 47$ (class) & density & ingredient & no & 491 & +0.225 [+0.189, +0.263] & +0.338 [+0.308, +0.366] & $-$0.113 [$-$0.148, $-$0.075] \\
WMC $\geq 47$ (class) & density & non-ingredient & yes & 491 & +0.178 [+0.148, +0.215] & +0.276 [+0.249, +0.301] & $-$0.098 [$-$0.126, $-$0.062] \\
PMD NPath (method) & size & non-ingredient & yes & 490 & $-$0.047 [$-$0.082, $-$0.007] & $-$0.050 [$-$0.058, $-$0.041] & +0.002 [$-$0.026, +0.039] \\
PMD NPath (method) & structure & non-ingredient & yes & 490 & $-$0.006 [$-$0.019, +0.010] & $-$0.024 [$-$0.033, $-$0.013] & +0.018 [+0.004, +0.034] \\
PMD NPath (method) & density & non-ingredient & yes & 490 & +0.109 [+0.086, +0.133] & +0.055 [+0.032, +0.080] & +0.054 [+0.041, +0.069] \\
Brain Method (method) & size & ingredient & no & 491 & +0.003 [$-$0.013, +0.028] & $-$0.014 [$-$0.025, $-$0.005] & +0.017 [0.000, +0.042] \\
Brain Method (method) & size & non-ingredient & yes & 491 & $-$0.030 [$-$0.049, $-$0.001] & $-$0.033 [$-$0.038, $-$0.028] & +0.003 [$-$0.014, +0.030] \\
Brain Method (method) & structure & ingredient & no & 491 & $-$0.040 [$-$0.058, $-$0.017] & $-$0.051 [$-$0.065, $-$0.034] & +0.011 [+0.002, +0.025] \\
Brain Method (method) & structure & non-ingredient & yes & 491 & +0.050 [+0.039, +0.061] & +0.011 [+0.005, +0.018] & +0.039 [+0.029, +0.052] \\
Brain Method (method) & density & ingredient & no & 491 & +0.047 [+0.013, +0.087] & +0.023 [$-$0.001, +0.046] & +0.024 [+0.005, +0.048] \\
Brain Method (method) & density & non-ingredient & yes & 491 & +0.110 [+0.093, +0.126] & +0.053 [+0.042, +0.064] & +0.056 [+0.046, +0.065] \\
\bottomrule
\end{tabular}
}
\end{table}

Table~\ref{tab:cells-human} gives the human cells and
Table~\ref{tab:cells-tool} the tool cells. The four predictions, read on
flow:

\begin{itemize}
\item \textbf{P1 \pOneHolds.} For human Blob, \dD{} is \pOneDPeak{} in the peak family and \pOneDDensity{} in density: the density family's
  lower bound is above 0, the peak family's is not, and the prediction asks
  for both. For \dO{} the prediction \pOneOHolds: \pOneOPeak{} in peak and
  \pOneODensity{} in density. Multiplying complexity by the data a class
  sends out raises its agreement with the reviewers in both families where a
  single complex method makes a class an outlier; multiplying it by the
  divergence does so in one, by a smaller margin.
\item \textbf{P2 \pTwoHolds.} \dD{} is larger in the peak family than in
  size for human Blob by \pTwoBlob{}, in structure than in size for Long
  Method by \pTwoLong{}, and in peak than in size for Fontana by
  \pTwoFontana{}; every tool cell's lower bound is above
  \pTwoMinToolLo{}. The factor is least harmful where the plan said it
  would help most, but that is because it damages the size families
  outright: in the size family \dD{} is \cellBlobSizeD{} for Blob,
  \cellLongSizeD{} for Long Method and \cellFontanaSizeD{} for Fontana,
  and its upper bound is below 0 in \cellsSizeDHiBelow{} of the
  \cellsSizeFlow{} counted size cells. On the JavaParser-only set the Blob
  contrast is \pTwoJpBlob{} and P2 \pTwoJpHolds{}.
\item \textbf{P3 reads \pThreeReads.} \dDO{} for human Blob is
  \pThreeSize{} in size, \pThreePeak{} in peak and \pThreeDensity{} in
  density: in every family the upper bound is below 0, so subtracting what
  comes in costs agreement rather than adding it. The competing expectation
  recorded from the exploratory observation, not the registered one, is
  what the data support.
\item \textbf{P4: \pFourPositive{} of \pFourCells{} counted tool cells
  positive, \pFourMajority.} The estimates split by family, not by oracle:
  the density cells are positive with lower bounds above 0 in every oracle
  and the size cells negative with upper bounds below 0 in every oracle
  (\pFourLoAbove{} cells resolved above 0, \pFourHiBelow{} below).
\end{itemize}

The verdict on the construct, by the pre-registered rule, is
\textbf{\consVerdict}: P1 holds for \dO{} and not for \dD{}, and P3 reads
outgoing-only. With constructors counted in the inputs the verdict is \consCtorVerdict{}
and \consCtorAgrees{}; on the JavaParser-only set it is \consJpVerdict{}
and \consJpAgrees{}. Section~\ref{sec:explore-forms} reads the \dO{} gain
threshold-free and at a matched flag rate, post hoc, and on MLCQ it does not
survive either. Under the data definition P1
\pOneDataHolds{} for \dD{} and \pOneDataOHolds{} for \dO{} (peak:
\pOneDataOPeak{}) and P3 reads \pThreeDataReads{}; under message, P1
\pOneMessageHolds{} for \dD{} and \pOneMessageOHolds{} for \dO{} (peak:
\pOneMessageOPeak{}) and P3 reads \pThreeMessageReads{}.

\paragraph{The pattern across cells.} Over the \cellsCountedFlow{} counted
cells under flow, \dDO{}'s upper bound is below 0 in \cellsCountedDOHiBelow{},
including every one of the \cellsCountedClassFlow{} class-level cells; its
lower bound is above 0 in \cellsCountedDOLoAbove{}, all at method level
against the tool and internal method definitions, where the divergence
removes false positives from a zero-heavy signal that flags almost
everything (Section~\ref{sec:tools}). \dO{}'s lower bound is above 0 in
\cellsCountedOLoAbove{} counted cells and its upper bound below 0 in
\cellsCountedOHiBelow{}; \dD{}'s in \cellsCountedDLoAbove{} and
\cellsCountedDHiBelow{}. The divergence factor's resolved gains,
\cellsCountedDLoAbove{} of the \cellsCountedFlow{} counted cells, lie
mostly in the density families (\cellsDensityDLoAbove{} of
\cellsDensityFlow{}), where a density signal alone is a poor god-unit
detector and any weighting that reintroduces size helps.
On Fontana's God Class, read as paired differences only, the outgoing
factor's gain is large in every family (density: \cellFontanaDensityO{})
and the divergence factor's smaller (\cellFontanaDensityD{}); in size the
divergence factor is \cellFontanaSizeD{}.

\paragraph{Between definitions.} Counting the call itself (flow against
data) and treating the first item as the message (message against flow)
change \dD{} by hundredths. Of \betweenCells{} contrasts of flow against
data, \betweenFlowDataLoAbove{} have a lower bound above 0 and
\betweenFlowDataHiBelow{} an upper bound below, all on the large seed-3
sample; of message against flow, \betweenMessageFlowLoAbove{} and
\betweenMessageFlowHiBelow{}. The human cells that resolve favour message
over flow by small margins, \betweenBlobPeakMessageFlow{} for Blob in
peak and \betweenFontanaDensityMessageFlow{} for Fontana in density. The
differences between definitions are an order of magnitude smaller than the
differences between the factors.

\subsection{Secondaries}
\label{sec:secondary}

\paragraph{Other definitions, variants and forms at class level.}
Every registered arm was read against \Xs{} on the primary's population;
Table~\ref{tab:other-arms} lists those not already in
Table~\ref{tab:primary}. Of the \secOtherArms{} arms, \secOtherLoAbove{}
have a pooled lower bound above 0, the outgoing factor under flow
(\secOFlowMicro{}) and under message (\secOMessageMicro{}), and
\secOtherHiBelow{} an upper bound below. The outgoing factor on \Xs{}
reaches an MCC of \secOFlowMcc{} against \primXMcc{}, recall
\secOFlowRecall{} against \primXRecall{} at precision \secOFlowPrecision{},
and its macro difference is \secOFlowMacro{}; its family average in
Section~\ref{sec:cons} is resolved as well. It flags \matchedOOwnFlags{}
reviewed classes where \Xs{} flags \matchedXOwnFlags{}; what that gain is
worth at a matched rate is read in Section~\ref{sec:explore-forms}. The divergence under data and
message, the public input variant, and both variants with constructors are
each within a few hundredths of the primary. The intersection of the two
outlier sets (a class flagged on \Xs{} and on $d$) raises precision and
halves recall: \secAndMicro{}. Form B, which aggregates the method-level
forms as the companion study aggregates complexity, has point estimates
between \secFormBMinMicro{} and \secFormBMaxMicro{} pooled over
\secFormBArms{} arms, none resolved, \secFormBMacroLoAbove{} of them with a
macro lower bound above 0; the per-file maximum of
\sig{nloc}$\cdot$\Dflow{} is \secDFormBMicro{} pooled and
\secDFormBMacro{} macro.

\begin{table}[t]
\centering
\caption{Every registered class-level arm against \Xs{} on the primary's
population, by \madz{}: the other two coupling definitions, the outgoing
factor, the factors alone, the intersection of the two outlier sets, the
public input variant, both variants with constructors, and class form B
(the per-file maximum of the method-level form). Secondaries: no
multiplicity correction, none can change the primary.}
\label{tab:other-arms}
\footnotesize
\setlength{\tabcolsep}{3pt}
\fitwidth{% generated by paper/scripts/tables.py; do not edit
\begin{tabular}{lrrlrrl}
\toprule
Form & MCC of X$^*$ & MCC of the form & Paired difference (95\% CI) & Precision & Recall & Reading \\
\midrule
nloc\_max$\cdot$D\_data (A, non-private) & 0.19 & 0.18 & $-$0.008 [$-$0.101, +0.096] & 0.20 & 0.47 & no difference resolved \\
nloc\_max$\cdot$D\_message (A, non-private) & 0.19 & 0.18 & $-$0.005 [$-$0.089, +0.084] & 0.20 & 0.48 & no difference resolved \\
nloc\_max$\cdot$O\_flow (A) & 0.19 & 0.27 & +0.085 [+0.001, +0.182] & 0.23 & 0.69 & the form refines complexity \\
nloc\_max$\cdot$O\_data (A) & 0.19 & 0.22 & +0.035 [$-$0.053, +0.134] & 0.21 & 0.61 & no difference resolved \\
nloc\_max$\cdot$O\_message (A) & 0.19 & 0.27 & +0.087 [+0.003, +0.183] & 0.23 & 0.68 & the form refines complexity \\
D\_flow (A, non-private) & 0.19 & 0.17 & $-$0.014 [$-$0.109, +0.091] & 0.21 & 0.40 & no difference resolved \\
D\_data (A, non-private) & 0.19 & 0.20 & +0.012 [$-$0.085, +0.124] & 0.24 & 0.38 & no difference resolved \\
D\_message (A, non-private) & 0.19 & 0.21 & +0.019 [$-$0.073, +0.123] & 0.23 & 0.43 & no difference resolved \\
O\_flow (A) & 0.19 & 0.19 & 0.000 [$-$0.100, +0.106] & 0.21 & 0.45 & no difference resolved \\
O\_data (A) & 0.19 & 0.21 & +0.021 [$-$0.069, +0.125] & 0.23 & 0.44 & no difference resolved \\
O\_message (A) & 0.19 & 0.22 & +0.034 [$-$0.064, +0.141] & 0.23 & 0.49 & no difference resolved \\
nloc\_max $\cap$ d\_flow (A, non-private) & 0.19 & 0.16 & $-$0.023 [$-$0.092, +0.045] & 0.30 & 0.17 & no difference resolved \\
nloc\_max $\cap$ d\_data (A, non-private) & 0.19 & 0.16 & $-$0.023 [$-$0.102, +0.057] & 0.33 & 0.15 & no difference resolved \\
nloc\_max $\cap$ d\_message (A, non-private) & 0.19 & 0.17 & $-$0.017 [$-$0.084, +0.048] & 0.31 & 0.18 & no difference resolved \\
nloc\_max$\cdot$D\_flow (A, public) & 0.19 & 0.16 & $-$0.027 [$-$0.116, +0.071] & 0.18 & 0.49 & no difference resolved \\
nloc\_max$\cdot$D\_data (A, public) & 0.19 & 0.19 & +0.003 [$-$0.092, +0.109] & 0.20 & 0.52 & no difference resolved \\
nloc\_max$\cdot$D\_message (A, public) & 0.19 & 0.20 & +0.015 [$-$0.070, +0.110] & 0.20 & 0.56 & no difference resolved \\
nloc\_max$\cdot$D\_flow (A, non-private, with constructors) & 0.19 & 0.12 & $-$0.062 [$-$0.130, +0.011] & 0.20 & 0.27 & no difference resolved \\
nloc\_max$\cdot$D\_data (A, non-private, with constructors) & 0.19 & 0.18 & $-$0.011 [$-$0.103, +0.090] & 0.23 & 0.36 & no difference resolved \\
nloc\_max$\cdot$D\_message (A, non-private, with constructors) & 0.19 & 0.14 & $-$0.047 [$-$0.125, +0.036] & 0.21 & 0.30 & no difference resolved \\
nloc\_max$\cdot$D\_flow (A, public, with constructors) & 0.19 & 0.15 & $-$0.038 [$-$0.128, +0.056] & 0.19 & 0.40 & no difference resolved \\
nloc\_max$\cdot$D\_data (A, public, with constructors) & 0.19 & 0.16 & $-$0.030 [$-$0.126, +0.077] & 0.20 & 0.40 & no difference resolved \\
nloc\_max$\cdot$D\_message (A, public, with constructors) & 0.19 & 0.17 & $-$0.020 [$-$0.109, +0.079] & 0.20 & 0.43 & no difference resolved \\
nloc$\cdot$D\_flow (B, non-private)\_max & 0.19 & 0.23 & +0.044 [$-$0.038, +0.139] & 0.21 & 0.60 & no difference resolved \\
nloc$\cdot$D\_data (B, non-private)\_max & 0.19 & 0.25 & +0.060 [$-$0.020, +0.154] & 0.23 & 0.59 & no difference resolved \\
nloc$\cdot$D\_message (B, non-private)\_max & 0.19 & 0.21 & +0.023 [$-$0.051, +0.105] & 0.21 & 0.55 & no difference resolved \\
nloc$\cdot$D\_flow (B, non-private, with constructors)\_max & 0.19 & 0.25 & +0.064 [$-$0.023, +0.162] & 0.22 & 0.63 & no difference resolved \\
nloc$\cdot$O\_flow (B)\_max & 0.19 & 0.21 & +0.026 [$-$0.046, +0.107] & 0.21 & 0.54 & no difference resolved \\
nloc$\cdot$O\_data (B)\_max & 0.19 & 0.18 & $-$0.003 [$-$0.075, +0.079] & 0.20 & 0.48 & no difference resolved \\
\bottomrule
\end{tabular}
}
\end{table}

\paragraph{Independence and strata.} Within confirm repositories \Xs{}
correlates with $d_{\text{flow}}$ at a median Spearman of \rhoXD{} and with
$\mathit{out}_{\text{flow}}$ at \rhoXOut{}; the number of methods
correlates with $d$ at \rhoNomD{} and with $\mathit{out}$ at \rhoNomOut{},
and the sum of lines with $\mathit{out}$ at \rhoNlocSumOut{}. The two
measures \trackClosely{} by the plan's criterion, and the primary re-read
within tertiles of \Xs{} per repository is \tertileOne{}, \tertileTwo{} and
\tertileThree{} from the least to the most complex third: the factor adds no
resolved information at any level of complexity. A class's $\mathit{out}$ is
a sum over its methods, so a large class sends more out whatever each method
does; that is the size the outgoing factor carries, and the size null
(Section~\ref{sec:primary}) shows how much of its effect that is.

\paragraph{Method level.} Against MLCQ Long Method on \methRepos{} confirm
repositories (\methUniv{} reviewed methods, \methPos{} positive), \Xm{} =
\sig{\xMethod} alone reaches an MCC of \methXMcc{} with precision
\methXPrecision{} and recall \methXRecall{}. Weighting it by \Dflow{} on all
calls gives \methMicro{}: \methReading{}, the point estimate against the
factor. The weighted form raises precision (\methDPrecision{}) and loses
recall (\methDRecall{}): it flags \methDTp{} of the \methPos{} Long Methods
against \methXTp{}, because a long method that calls nothing in its own
repository is a catchment and is exempt. The no-helpers variant is
\methNoHelpMicro{}, \methNoHelpReading{}; the two variants differ by
\methNoHelpVsAllMicro{}. The halves say which part costs: exemption only is
\methEMicro{} (MCC \methEMcc{}, recall \methERecall{}), within hundredths
of complexity alone; amplification only is \methAmpMicro{},
\methAmpReading{}; the whole factor against the exemption alone is
\methDvsEMicro{}, \methDvsEReading{}, and against the amplification alone
\methDvsAMicro{}. At method level the damage is the amplification, which
hands the flag to methods that call much and are not long, and the support
the leaf diagnostic gives the exemption (Table~\ref{tab:leaf}) does not
reach it. Fitted on the non-exempt units only, with the exempt never
flagged, the weighted form recovers recall (\methNeRecall{}) at a cost in
precision (\methNePrecision{}): \methNeVsXMicro{} against complexity alone,
\methNeVsXReading{}. The mechanism is not read, since the primary does not
read refinement. The JavaParser-only set agrees: \methJpMicro{},
\methJpReading{}.

\begin{table}[t]
\centering
\caption{Method level against MLCQ human Long Method on the confirm half,
\methUniv{} reviewed methods in \methRepos{} repositories: \Xm{} alone and
its forms under flow on all calls, by \madz{}.}
\label{tab:method}
\small
\setlength{\tabcolsep}{4pt}
\fitwidth{% generated by paper/scripts/tables.py; do not edit
\begin{tabular}{lrrrrl}
\toprule
Arm (method level, flow, all calls) & MCC & Precision & Recall & Flagged & Paired difference against X$^*_m$ (95\% CI) \\
\midrule
X$^*_m$ = \texttt{cognitive}, complexity alone & 0.603 & 0.69 & 0.61 & 110 &  \\
X$^*_m\cdot$D$_{\text{flow}}$, all calls & 0.511 & 0.86 & 0.34 & 50 & $-$0.092 [$-$0.188, +0.010], no difference resolved \\
X$^*_m\cdot$D$_{\text{flow}}$, no helpers & 0.460 & 0.76 & 0.33 & 54 & $-$0.144 [$-$0.244, $-$0.043], it makes it worse \\
X$^*_m\cdot$E$_{\text{flow}}$, exemption only & 0.601 & 0.72 & 0.58 & 100 & $-$0.002 [$-$0.077, +0.073], no difference resolved \\
X$^*_m\cdot$A$_{\text{flow}}$, amplification only & 0.510 & 0.76 & 0.40 & 66 & $-$0.093 [$-$0.174, $-$0.007], the amplification makes it worse \\
X$^*_m\cdot$D$_{\text{flow}}$, fitted on the non-exempt units & 0.543 & 0.54 & 0.68 & 158 & $-$0.060 [$-$0.153, +0.045], no difference resolved \\
\bottomrule
\end{tabular}
}
\end{table}

\paragraph{Tool oracles.}
\label{sec:tools}
Table~\ref{tab:tool-oracles} reads the primary's paired difference on the
seed-3 sample against each tool definition, never the headline and read with
the caveat that PMD's GodClass contains a coupling measure. At class level
the divergence factor makes every one worse by a resolved margin: PMD
GodClass \toolPmdMicro{} over \toolPmdN{} repositories, DECOR Blob
\toolDecorMicro{}, WMC $\geq 47$ \toolWmcMicro{}; with constructors counted
the losses are smaller and still resolved (\toolPmdCtorMicro{},
\toolDecorCtorMicro{}, \toolWmcCtorMicro{}). In each the factor raises
recall and halves precision, as against the human labels. At method level
\Xm{} alone is a recall-heavy flag (PMD NPath: recall \toolNpathXRecall{}
at precision \toolNpathXPrecision{}; Brain Method: recall
\toolBrainXRecall{} at precision \toolBrainXPrecision{}), because cognitive
complexity ties at zero for most methods and the rule's fallback flags a
large tail. Against NPath the factor is \toolNpathMicro{}, not resolved;
against Brain Method, whose definition contains nesting, lines and
cyclomatic complexity, \toolBrainMicro{}, \toolBrainReading{} on both sets
(\toolBrainJpMicro{} on JavaParser-only): the one resolved gain of the
primary arm against any oracle, on an internal definition, and it is the exemption
pruning false positives from a flag set of \toolBrainXRecall{} recall, not a
better ranking of god methods.

\begin{table}[t]
\centering
\caption{The primary form against the tool and internal definitions on the
seed-3 sample, by \madz{}, pooled over the repositories carrying each
definition. Positives are the definition's own flags. PMD GodClass contains
a coupling measure (ATFD), so agreement with it is partly circular.}
\label{tab:tool-oracles}
\small
\setlength{\tabcolsep}{4pt}
\fitwidth{% generated by paper/scripts/tables.py; do not edit
\begin{tabular}{lrrrrlll}
\toprule
Oracle (seed-3 sample) & $n$ & Positives & MCC of X$^*$ & MCC of X$^*\cdot$D$_{\text{flow}}$ & Paired difference (95\% CI) & Reading & With constructors \\
\midrule
PMD GodClass (class) & 490 & 18{,}710 & 0.351 & 0.208 & $-$0.143 [$-$0.176, $-$0.106] & it makes it worse & $-$0.117 [$-$0.153, $-$0.078] \\
DECOR Blob (class) & 490 & 26{,}372 & 0.202 & 0.109 & $-$0.093 [$-$0.121, $-$0.059] & it makes it worse & $-$0.051 [$-$0.097, $-$0.002] \\
WMC $\geq 47$ (class) & 491 & 31{,}633 & 0.390 & 0.247 & $-$0.144 [$-$0.185, $-$0.099] & it makes it worse & $-$0.122 [$-$0.161, $-$0.074] \\
PMD NPath (method) & 490 & 23{,}707 & 0.350 & 0.330 & $-$0.021 [$-$0.054, +0.013] & no difference resolved & -- \\
Brain Method (method) & 491 & 9{,}440 & 0.220 & 0.256 & +0.036 [+0.025, +0.047] & the divergence factor refines complexity & -- \\
\bottomrule
\end{tabular}
}
\end{table}

\paragraph{Fontana's God Class.} On the \fontanaPopSystems{} systems
carrying a labelled class (\fontanaPopUniv{} reviewed classes,
\fontanaPopPos{} God Classes), \Xs{} alone reaches an MCC of \fontanaXMcc{}, precision
\fontanaXPrecision{} and recall \fontanaXRecall{}, the balanced and extreme
instances separating as expected. The primary's paired difference is
\fontanaMicro{}: \fontanaReading{}; the weighted arm's MCC is
\fontanaDMcc{}, its false positives \fontanaDFp{} against \fontanaXFp{}.
With constructors, \fontanaCtorMicro{}. The leaf diagnostic
(Table~\ref{tab:leaf}) is \leafFontanaReading{} under flow. The replication on a different team's
labels, different systems and older code gives the same shape as MLCQ: a
negative point estimate on the whole factor, a large resolved gain on the
outgoing factor (\cellFontanaDensityO{} in the density family), and a
resolved loss when what comes in is subtracted (\cellFontanaDensityDO{}).

\paragraph{The matrix by family.} Table~\ref{tab:friedman} gives, for each
oracle and family under flow, a Friedman test over repositories across the
three forms (X, X$\cdot$D, X$\cdot$O), signal by signal within the family
\cite{friedman1937}. Against human Blob the forms are not separated in the
size family (\friedBlobSizeP{}) and are in peak (\friedBlobPeakP{})
and density (\friedBlobDensityP{}), where X$\cdot$O has the best mean
rank (\friedBlobPeakRankO{} against \friedBlobPeakRankX{} for X and
\friedBlobPeakRankD{} for X$\cdot$D in peak; a lower rank is better).
Against Long Method the size family separates with X best
(\friedLongSizeRankX{} against \friedLongSizeRankD{} for X$\cdot$D), the
structure family does not separate (\friedLongStructureP{}), and the
density family separates with both weighted forms ahead of X. On Fontana
the peak and density families separate with X$\cdot$O best
(\friedFontanaDensityRankO{} against \friedFontanaDensityRankX{}). A claim
that the factor helps in general would have to hold within a family that is
not size; what holds there is that the outgoing factor helps and the
divergence factor sits between it and complexity alone.

\begin{table}[t]
\centering
\caption{Friedman tests over repositories across the three forms within
each family, flow definition, \madz{}: mean rank of each form (1 is best),
$n$ repositories after dropping those where every signal ties, and the
$\chi^2$ test's $p$.}
\label{tab:friedman}
\small
\setlength{\tabcolsep}{4pt}
\fitwidth{% generated by paper/scripts/tables.py; do not edit
\begin{tabular}{llrrrrr}
\toprule
Oracle & Family & $n$ & Mean rank X & Mean rank X$\cdot$D & Mean rank X$\cdot$O & $p$ (Friedman $\chi^2$) \\
\midrule
MLCQ Blob (class) & size & 53 & 2.08 & 2.16 & 1.76 & 0.100 \\
MLCQ Blob (class) & peak & 52 & 2.32 & 2.11 & 1.58 & $<$0.001 \\
MLCQ Blob (class) & density & 53 & 2.37 & 2.00 & 1.63 & $<$0.001 \\
MLCQ Long Method (method) & size & 63 & 1.40 & 2.40 & 2.20 & $<$0.001 \\
MLCQ Long Method (method) & structure & 65 & 1.95 & 2.14 & 1.92 & 0.387 \\
MLCQ Long Method (method) & density & 66 & 2.55 & 1.71 & 1.74 & $<$0.001 \\
Fontana God Class (class) & size & 44 & 1.99 & 2.09 & 1.92 & 0.723 \\
Fontana God Class (class) & peak & 48 & 2.65 & 1.85 & 1.50 & $<$0.001 \\
Fontana God Class (class) & density & 48 & 2.84 & 1.85 & 1.30 & $<$0.001 \\
PMD GodClass (class) & size & 387 & 1.21 & 2.64 & 2.15 & $<$0.001 \\
PMD GodClass (class) & peak & 387 & 2.16 & 2.33 & 1.51 & $<$0.001 \\
PMD GodClass (class) & density & 387 & 2.61 & 2.09 & 1.30 & $<$0.001 \\
DECOR Blob (class) & size & 439 & 1.23 & 2.78 & 1.99 & $<$0.001 \\
DECOR Blob (class) & peak & 439 & 2.17 & 2.37 & 1.47 & $<$0.001 \\
DECOR Blob (class) & density & 439 & 2.60 & 2.04 & 1.37 & $<$0.001 \\
WMC $\geq 47$ (class) & size & 421 & 1.20 & 2.68 & 2.12 & $<$0.001 \\
WMC $\geq 47$ (class) & peak & 421 & 2.20 & 2.29 & 1.51 & $<$0.001 \\
WMC $\geq 47$ (class) & density & 421 & 2.66 & 2.00 & 1.34 & $<$0.001 \\
PMD NPath (method) & size & 363 & 1.69 & 1.73 & 2.59 & $<$0.001 \\
PMD NPath (method) & structure & 363 & 1.91 & 1.78 & 2.31 & $<$0.001 \\
PMD NPath (method) & density & 363 & 2.36 & 1.44 & 2.20 & $<$0.001 \\
Brain Method (method) & size & 271 & 1.75 & 1.57 & 2.68 & $<$0.001 \\
Brain Method (method) & structure & 271 & 1.97 & 1.68 & 2.34 & $<$0.001 \\
Brain Method (method) & density & 271 & 2.33 & 1.41 & 2.26 & $<$0.001 \\
\bottomrule
\end{tabular}
}
\end{table}

\subsection{Sensitivities}
\label{sec:sens}

\paragraph{Dependency injection.} Without the confirm repositories that
declare a DI framework, \sensDiRepos{} repositories and \sensDiUniv{}
reviewed classes remain, and the primary is \sensDiMicro{},
\sensDiReading{}, with the arms at \sensDiXMcc{} and \sensDiDMcc{}. The
point estimate moves further against the factor, not towards it, when the
calls static parsing cannot see are removed with their repositories.

\paragraph{The scale.} Fitting the weighted arm on the non-exempt units
only, with the exempt units never flagged, gives an arm at MCC
\sensNeMcc{} (precision \sensNePrecision{}, recall \sensNeRecall{},
\sensNeFp{} false positives): \sensNeVsXMicro{} against \Xs{},
\sensNeVsXReading{}, and \sensNeVsAllMicro{} against the same arm fitted
on every unit, \sensNeVsAllReading{} (\sensJpNeVsAllMicro{} on the
JavaParser-only set). The mean-absolute-deviation fallback that the
zeros force is part of the primary arm's loss, and removing it brings the
arm back to complexity alone, not past it.

\paragraph{Coupled units only and the time budget.} Every reviewed confirm
class carries coupling, so the primary on the coupled classes only is the
primary (\sensUncoupledMicro{}); no repository spent its resolution budget,
so the primary without such repositories is the primary.

\subsection{Exploratory: the intersection of the two outlier sets}
\label{sec:inter}

Table~\ref{tab:intersection} reads the positives by region when complexity
and divergence are each put through the rule on their own. Against human
Blob the \interBlobBothTp{} hits among the \interBlobBothTp{} +
\interBlobBothFp{} classes flagged on both are \interBlobBothPrecision{}
precise, against \interBlobXOnlyPrecision{} for the classes flagged on
complexity only and \interBlobDOnlyPrecision{} on divergence only; the
divergence alone, as a detector, has MCC \interBlobDMcc{} against
\primXMcc{} for complexity alone (\interBlobDVsX{}). The union is
\interBlobEitherVsX{} against complexity alone and the intersection
\interBlobBothVsX{}, neither resolved. Against Long Method the intersection
is \interLongBothPrecision{} precise, and the divergence alone is a poorer
detector than complexity alone (MCC \interLongDMcc{}, \interLongDVsX{});
on Fontana the union gains (\interFontanaEitherVsX{}, MCC
\interFontanaEitherMcc{}) and the intersection, \interFontanaBothPrecision{}
precise, loses recall (\interFontanaBothVsX{}). The units in both outlier
sets are the most often god units under every oracle, by margins the human
samples do not resolve. This is exploration, labelled as such, with no
prediction and no verdict.

\begin{table}[t]
\centering
\caption{Exploratory: the reviewed positives by region when the outliers of
\Xs{} (or \Xm{}) and of the second set are flagged separately by \madz{}.
Each region is hits/false positives among the reviewed units; the last
column is the pooled MCC of the union minus that of \Xs{}.}
\label{tab:intersection}
\small
\setlength{\tabcolsep}{4pt}
\fitwidth{% generated by paper/scripts/tables.py; do not edit
\begin{tabular}{llrrrrrrl}
\toprule
Oracle & Second set & X only & $d$ only & Both & Either & MCC both & MCC either & Either against X (95\% CI) \\
\midrule
MLCQ Blob & $d_{\text{flow}}$ & 13/48 & 15/54 & 17/39 & 45/141 & 0.164 & 0.223 & +0.036 [$-$0.023, +0.106] \\
MLCQ Blob & O$_{\text{flow}}$ & 11/37 & 25/113 & 19/50 & 55/200 & 0.157 & 0.221 & +0.034 [$-$0.041, +0.118] \\
MLCQ Long Method & $d_{\text{flow}}$ & 29/27 & 30/80 & 47/7 & 106/114 & 0.540 & 0.576 & $-$0.028 [$-$0.091, +0.043] \\
MLCQ Long Method & O$_{\text{flow}}$ & 28/23 & 25/79 & 48/11 & 101/113 & 0.523 & 0.550 & $-$0.053 [$-$0.117, +0.018] \\
Fontana God Class & $d_{\text{flow}}$ & 24/10 & 12/3 & 85/1 & 121/14 & 0.699 & 0.818 & +0.052 [+0.011, +0.094] \\
Fontana God Class & O$_{\text{flow}}$ & 10/9 & 20/11 & 99/2 & 129/22 & 0.768 & 0.824 & +0.058 [$-$0.001, +0.119] \\
\bottomrule
\end{tabular}
}
\end{table}

\subsection{Exploratory, post hoc: size times calls, and the leaf gate}
\label{sec:explore-forms}

Everything in this subsection was computed after every oracle of the study
had been read, from the raw records, by a script that is in the ancillary
data with its output and digest; it is the exploratory step of a next
registration and changes nothing above. It asks what the registered forms
mixed together: whether the signal in the outgoing factor is the calls or
the data on them, whether it is beyond size, and what the exemption does on
its own. Each form is scored by the registered rule on every unit of the
repository and read against the best complexity-alone signal and against a
size-only counterpart, with the plan's bootstrap, on both MLCQ halves (both
exploration material now), on Fontana's God Class and on the tool oracles.

\begin{table}[t]
\centering
\caption{Exploratory, post hoc: class forms built from size and outgoing
calls, pooled MCC against MLCQ human Blob on each half and against Fontana's
God Class, by \madz{}. ``Leaving calls'' is the number of call sites in the
class that resolve to the repository's own code outside the class, each call
one unit; ``fan-out'' the number of distinct callee files. The paired
differences are against Halstead volume summed (the best complexity-alone
signal) and against lines summed $\times$ number of methods (size only), on
the confirm half.}
\label{tab:explore-class}
\footnotesize
\setlength{\tabcolsep}{3pt}
\fitwidth{% generated by paper/scripts/tables.py; do not edit
\begin{tabular}{lrrllr}
\toprule
Class form & Explore MCC & Confirm MCC & Confirm, against complexity alone & Confirm, against size only & Fontana MCC \\
\midrule
Halstead volume, summed (complexity alone) & 0.287 & 0.254 & -- & +0.037 [$-$0.015, +0.095] & 0.835 \\
lines, summed & 0.266 & 0.232 & $-$0.022 [$-$0.068, +0.015] & +0.015 [$-$0.030, +0.059] & 0.827 \\
cyclomatic complexity, summed (WMC) & 0.195 & 0.233 & $-$0.021 [$-$0.069, +0.024] & +0.016 [$-$0.040, +0.066] & 0.812 \\
lines, per-file maximum (X$^*$) & 0.287 & 0.187 & $-$0.067 [$-$0.134, $-$0.001] & $-$0.030 [$-$0.109, +0.049] & 0.765 \\
cognitive complexity, maximum & 0.234 & 0.161 & $-$0.093 [$-$0.169, $-$0.017] & $-$0.056 [$-$0.139, +0.028] & 0.810 \\
cyclomatic density, maximum & 0.075 & 0.039 & $-$0.215 [$-$0.327, $-$0.115] & $-$0.178 [$-$0.281, $-$0.091] & 0.202 \\
number of methods & 0.154 & 0.160 & $-$0.094 [$-$0.182, $-$0.019] & $-$0.057 [$-$0.126, +0.007] & 0.590 \\
leaving calls alone & 0.237 & 0.214 & $-$0.040 [$-$0.111, +0.034] & $-$0.003 [$-$0.069, +0.071] & 0.812 \\
fan-out (callee files) alone & 0.218 & 0.158 & $-$0.096 [$-$0.175, $-$0.019] & $-$0.058 [$-$0.128, +0.015] & 0.847 \\
lines summed $\times$ number of methods (size only) & 0.286 & 0.217 & $-$0.037 [$-$0.095, +0.015] & -- & 0.720 \\
lines maximum $\times$ number of methods (X$^*\cdot$N) & 0.281 & 0.236 & $-$0.018 [$-$0.069, +0.031] & +0.020 [$-$0.015, +0.054] & 0.776 \\
lines summed $\times$ (1 + leaving calls) & 0.314 & 0.286 & +0.032 [$-$0.037, +0.102] & +0.069 [+0.010, +0.135] & 0.822 \\
lines summed $\times$ (1 + fan-out) & 0.317 & 0.280 & +0.026 [$-$0.041, +0.093] & +0.063 [+0.008, +0.124] & 0.827 \\
lines summed $\times$ (1 + all internal calls) & 0.325 & 0.269 & +0.015 [$-$0.045, +0.076] & +0.052 [$-$0.008, +0.117] & 0.800 \\
lines maximum $\times$ (1 + leaving calls) & 0.293 & 0.275 & +0.021 [$-$0.048, +0.098] & +0.058 [$-$0.011, +0.136] & 0.815 \\
X$^*\cdot$O$_{\text{flow}}$ (registered, argument-weighted) & 0.296 & 0.272 & +0.018 [$-$0.055, +0.094] & +0.056 [$-$0.017, +0.136] & 0.837 \\
cyclomatic summed $\times$ (1 + leaving calls) & 0.322 & 0.266 & +0.012 [$-$0.059, +0.080] & +0.049 [$-$0.006, +0.109] & 0.820 \\
cyclomatic density maximum $\times$ (1 + leaving calls) & 0.242 & 0.229 & $-$0.025 [$-$0.106, +0.053] & +0.013 [$-$0.055, +0.085] & 0.830 \\
number of methods $\times$ (1 + leaving calls) & 0.298 & 0.275 & +0.021 [$-$0.046, +0.085] & +0.058 [+0.004, +0.115] & 0.789 \\
lines maximum $\times$ (1 + leaving calls per method) & 0.298 & 0.193 & $-$0.061 [$-$0.136, +0.019] & $-$0.024 [$-$0.102, +0.059] & 0.753 \\
sum over methods of lines $\times$ (1 + leaving calls), form B & 0.262 & 0.241 & $-$0.013 [$-$0.072, +0.046] & +0.024 [$-$0.036, +0.092] & 0.827 \\
largest over methods of lines $\times$ (1 + leaving calls), form B & 0.263 & 0.188 & $-$0.066 [$-$0.143, +0.014] & $-$0.029 [$-$0.103, +0.051] & 0.804 \\
\bottomrule
\end{tabular}
}
\end{table}

\paragraph{Classes.} Of \expClassForms{} forms (Table~\ref{tab:explore-class}),
the one that leads on both halves is the class's total lines times one plus
its leaving calls: \expSizeCallsExploreMcc{} on explore and
\expSizeCallsConfirmMcc{} on confirm, against \expComplexityExploreMcc{}
and \expComplexityConfirmMcc{} for the best complexity-alone signal
(\expSizeCallsConfirmVsComplexity{} on confirm, \expSizeCallsExploreVsComplexity{}
on explore, neither resolved) and against \expSizeOnlyExploreMcc{} and
\expSizeOnlyConfirmMcc{} for the size-only counterpart
(\expSizeCallsConfirmVsSize{} on confirm, resolved;
\expSizeCallsExploreVsSize{} on explore, not). The plain count of calls does
as well as the registered argument-weighted outgoing data, and fan-out by
distinct callee files does as well as the call count
(\expSizeFanConfirmMcc{}); so the unit of the signal is the call, not the
data on it. The calls alone are weak (\expCallsAloneConfirmMcc{}), the
per-method products summed (form B) are below the class total
(\expFormBSumConfirmMcc{}), and the per-method call rate, which takes size
out, is weaker still (\expRateConfirmMcc{}). Cyclomatic density times calls
(\expDensityCallsConfirmMcc{}) says the complexity beyond length adds
nothing. The registered \Xs$\cdot$\Oflow{} against its own size null
\Xs$\cdot$N, the comparison the plan did not register, is
\expXOConfirmVsSizeMax{} on confirm and \expXOExploreVsSizeMax{} on explore,
and \expXOFontanaVsSizeMax{} on Fontana. On Fontana every form of size ties
near \expFontanaComplexityMcc{}. Against the tool oracles the ingredient
signals lead, as they must: \expToolPmdTop{} against PMD GodClass
(\expToolPmdTopMcc{}), \expToolWmcTop{} against WMC (\expToolWmcTopMcc{}).

\begin{table}[t]
\centering
\caption{Exploratory, post hoc: method forms against MLCQ human Long Method
on each half, by \madz{}. A gate ``[makes a call]'' multiplies by 1 or 0;
``internal'' means a call that resolves to the repository's own code,
``any'' any call site whatever its target. Precision and recall are on the
confirm half; the paired differences are against tokens (the best
complexity-alone signal) and against lines.}
\label{tab:explore-method}
\footnotesize
\setlength{\tabcolsep}{3pt}
\fitwidth{% generated by paper/scripts/tables.py; do not edit
\begin{tabular}{lrrrrll}
\toprule
Method form & Explore MCC & Confirm MCC & Precision & Recall & Confirm, against tokens & Confirm, against lines \\
\midrule
tokens (complexity alone) & 0.620 & 0.626 & 0.48 & 0.97 & -- & +0.105 [+0.048, +0.162] \\
cyclomatic complexity & 0.531 & 0.612 & 0.69 & 0.62 & $-$0.014 [$-$0.109, +0.075] & +0.092 [$-$0.003, +0.181] \\
cognitive complexity (X$^*_m$) & 0.600 & 0.603 & 0.69 & 0.61 & $-$0.022 [$-$0.125, +0.071] & +0.083 [$-$0.011, +0.171] \\
lines & 0.477 & 0.520 & 0.37 & 0.98 & $-$0.105 [$-$0.161, $-$0.048] & -- \\
Halstead volume & 0.573 & 0.583 & 0.43 & 0.98 & $-$0.043 [$-$0.063, $-$0.021] & +0.063 [+0.008, +0.114] \\
cyclomatic density & $-$0.044 & $-$0.042 & 0.04 & 0.01 & $-$0.668 [$-$0.737, $-$0.590] & $-$0.562 [$-$0.625, $-$0.493] \\
lines $\times$ [makes any call] & 0.679 & 0.744 & 0.69 & 0.87 & +0.119 [+0.070, +0.171] & +0.224 [+0.164, +0.280] \\
lines $\times$ [makes an internal call] & 0.679 & 0.697 & 0.70 & 0.78 & +0.072 [+0.009, +0.142] & +0.177 [+0.109, +0.250] \\
lines $\times$ [makes a call leaving the class] & 0.601 & 0.561 & 0.64 & 0.58 & $-$0.065 [$-$0.145, +0.018] & +0.041 [$-$0.043, +0.122] \\
tokens $\times$ [makes an internal call] & 0.581 & 0.572 & 0.52 & 0.77 & $-$0.054 [$-$0.123, +0.023] & +0.052 [$-$0.018, +0.131] \\
cyclomatic $\times$ [makes an internal call] & 0.579 & 0.613 & 0.73 & 0.59 & $-$0.012 [$-$0.099, +0.076] & +0.093 [+0.004, +0.188] \\
cognitive $\times$ [makes an internal call] & 0.548 & 0.610 & 0.72 & 0.59 & $-$0.016 [$-$0.101, +0.071] & +0.090 [+0.003, +0.172] \\
cyclomatic density $\times$ [makes an internal call] & 0.058 & 0.017 & 0.15 & 0.04 & $-$0.609 [$-$0.696, $-$0.509] & $-$0.503 [$-$0.583, $-$0.416] \\
lines $\times$ (1 + internal calls), the weight & 0.459 & 0.465 & 0.33 & 0.94 & $-$0.161 [$-$0.215, $-$0.110] & $-$0.055 [$-$0.105, $-$0.010] \\
internal calls alone & 0.411 & 0.488 & 0.55 & 0.55 & $-$0.138 [$-$0.230, $-$0.044] & $-$0.033 [$-$0.119, +0.059] \\
\bottomrule
\end{tabular}
}
\end{table}

\paragraph{Methods.} The method-level result is the exemption in its
crudest form, and it is the largest effect in the paper
(Table~\ref{tab:explore-method}). Among methods that make at least one call
of any kind, a method's length alone reaches an MCC of
\expGateAnyConfirmMcc{} on confirm (precision \expGateAnyConfirmPrecision{},
recall \expGateAnyConfirmRecall{}) and \expGateAnyExploreMcc{} on explore,
against \expTokensConfirmMcc{} and \expTokensExploreMcc{} for tokens, the
best complexity-alone signal: \expGateAnyConfirmVsTokens{} on confirm and
\expGateAnyExploreVsTokens{} on explore. Gating on an internal call instead
gives \expGateIntConfirmMcc{} and \expGateIntExploreMcc{}
(\expGateIntConfirmVsTokens{} and \expGateIntExploreVsTokens{}), so the gain
does not depend on the symbol solver; gating on a call that leaves the class
loses it (\expGateLeaveConfirmMcc{}). Weighting lines by the number of calls,
the registered form, destroys it: \expWeightConfirmMcc{},
\expWeightConfirmVsTokens{}. Lines alone, ungated, is \expLinesConfirmMcc{}.
Cyclomatic density is near zero gated or not (\expDensityGateConfirmMcc{},
\expMethodDensityConfirmMcc{}). What the gate does is not what it looks
like. Methods that make no call at all are almost never long
(\revNoCallLong{} long ones among \revNoCallReviews{} reviews of such
methods), so the gate removes nothing the reviewers would have called a
Long Method; what it removes is the short no-call methods, the getters,
setters and trivial constructors, whose zeros move the repository's median
and MAD so that the constant lands where length separates best. Raw lines
at 3.5 flag \matchedLinesOwnFlagsConfirm{} reviewed methods, tokens
\matchedTokensOwnFlagsConfirm{}, the gate \matchedGateOwnFlagsConfirm{};
the gate's gain is a recalibration of the constant, and the next paragraph
shows it ranks no better than length. Against the tool definitions, which have no
such exemption, the gate costs a little (PMD NPath: \expToolNpathCcMcc{} to
\expToolNpathCcGateMcc{} on cyclomatic; Brain Method:
\expToolBrainCogMcc{} to \expToolBrainCogGateMcc{} on cognitive).

\begin{table}[t]
\centering
\caption{Exploratory, post hoc: ranking or flag rate, class level against
MLCQ human Blob. AUC is the threshold-free pooled AUC of the form's
per-repository modified z over the reviewed units, with the paired bootstrap
difference against \Xs{} on the confirm half. ``Matched'' reads every form
at the constant that gives it as many flags over every class of the confirm
repositories as \Xs$\cdot$\Oflow{} has at 3.5 (\matchedClassFlags{} of
\matchedClassUnits{}, \matchedClassRate{}).}
\label{tab:matched-class}
\footnotesize
\setlength{\tabcolsep}{3pt}
\fitwidth{% generated by paper/scripts/tables.py; do not edit
\begin{tabular}{lrlrrrrl}
\toprule
Form & AUC, explore & AUC, confirm (vs X$^*$) & Flags at 3.5 & MCC at 3.5 & MCC, matched & Recall, matched & MCC vs X$^*$, matched (95\% CI) \\
\midrule
X$^*$ = \texttt{nloc\_max} & 0.804 & 0.768 (--) & 117 & 0.187 & 0.240 & 0.64 & -- \\
X$^*\cdot$D$_{\text{flow}}$, the primary arm & 0.651 & 0.631 ($-$0.137 [$-$0.195, $-$0.078]) & 214 & 0.146 & 0.162 & 0.47 & $-$0.078 [$-$0.150, $-$0.003] \\
X$^*\cdot$E$_{\text{flow}}$ & 0.669 & 0.640 ($-$0.128 [$-$0.183, $-$0.069]) & 138 & 0.212 & 0.191 & 0.52 & $-$0.049 [$-$0.124, +0.028] \\
X$^*\cdot$A$_{\text{flow}}$ & 0.761 & 0.740 ($-$0.029 [$-$0.062, +0.009]) & 279 & 0.196 & 0.203 & 0.56 & $-$0.036 [$-$0.114, +0.042] \\
X$^*\cdot$N, the size null & 0.795 & 0.770 (+0.002 [$-$0.027, +0.035]) & 270 & 0.236 & 0.292 & 0.77 & +0.052 [$-$0.005, +0.110] \\
X$^*\cdot$O$_{\text{flow}}$, the outgoing factor & 0.774 & 0.763 ($-$0.005 [$-$0.043, +0.033]) & 301 & 0.272 & 0.272 & 0.69 & +0.032 [$-$0.036, +0.099] \\
Halstead volume summed & 0.812 & 0.783 (+0.015 [$-$0.010, +0.046]) & 214 & 0.254 & 0.277 & 0.77 & +0.037 [$-$0.023, +0.100] \\
lines summed & 0.805 & 0.775 (+0.007 [$-$0.022, +0.043]) & 205 & 0.232 & 0.259 & 0.73 & +0.019 [$-$0.041, +0.085] \\
lines summed $\times$ number of methods & 0.772 & 0.746 ($-$0.022 [$-$0.061, +0.023]) & 320 & 0.217 & 0.226 & 0.69 & $-$0.014 [$-$0.085, +0.053] \\
lines summed $\times$ (1 + leaving calls) & 0.792 & 0.776 (+0.008 [$-$0.028, +0.048]) & 332 & 0.286 & 0.286 & 0.76 & +0.046 [$-$0.018, +0.111] \\
lines summed $\times$ (1 + fan-out) & 0.795 & 0.774 (+0.006 [$-$0.030, +0.050]) & 308 & 0.280 & 0.255 & 0.72 & +0.015 [$-$0.048, +0.077] \\
leaving calls alone & 0.721 & 0.721 ($-$0.047 [$-$0.102, +0.008]) & 206 & 0.214 & 0.216 & 0.62 & $-$0.024 [$-$0.104, +0.057] \\
\bottomrule
\end{tabular}
}
\end{table}

\begin{table}[t]
\centering
\caption{The same at method level against MLCQ human Long Method, matched to
the flag rate of tokens at 3.5 (\matchedMethodRate{} of methods), against
\Xm{}.}
\label{tab:matched-method}
\footnotesize
\setlength{\tabcolsep}{3pt}
\fitwidth{% generated by paper/scripts/tables.py; do not edit
\begin{tabular}{lrlrrrrl}
\toprule
Form & AUC, explore & AUC, confirm (vs X$^*$) & Flags at 3.5 & MCC at 3.5 & MCC, matched & Recall, matched & MCC vs X$^*$, matched (95\% CI) \\
\midrule
X$^*_m$ = \texttt{cognitive} & 0.894 & 0.936 (--) & 110 & 0.603 & 0.553 & 0.90 & -- \\
X$^*_m\cdot$D$_{\text{flow}}$, the primary arm & 0.790 & 0.832 ($-$0.105 [$-$0.158, $-$0.056]) & 50 & 0.511 & 0.449 & 0.76 & $-$0.104 [$-$0.180, $-$0.033] \\
X$^*_m\cdot$E$_{\text{flow}}$ & 0.787 & 0.829 ($-$0.107 [$-$0.159, $-$0.059]) & 100 & 0.601 & 0.428 & 0.75 & $-$0.125 [$-$0.197, $-$0.058] \\
X$^*_m\cdot$O$_{\text{flow}}$ & 0.880 & 0.915 ($-$0.021 [$-$0.036, $-$0.005]) & 76 & 0.555 & 0.500 & 0.84 & $-$0.053 [$-$0.106, $-$0.001] \\
tokens & 0.962 & 0.963 (+0.027 [$-$0.000, +0.054]) & 252 & 0.626 & 0.626 & 0.97 & +0.072 [+0.009, +0.133] \\
lines & 0.931 & 0.956 (+0.020 [$-$0.005, +0.045]) & 332 & 0.520 & 0.565 & 0.91 & +0.011 [$-$0.049, +0.072] \\
lines $\times$ [makes any call] & 0.953 & 0.975 (+0.039 [+0.013, +0.063]) & 157 & 0.744 & 0.607 & 0.98 & +0.054 [$-$0.002, +0.105] \\
lines $\times$ [makes an internal call] & 0.855 & 0.855 ($-$0.081 [$-$0.145, $-$0.025]) & 139 & 0.697 & 0.514 & 0.83 & $-$0.039 [$-$0.119, +0.032] \\
lines $\times$ (1 + internal calls) & 0.931 & 0.936 (+0.000 [$-$0.033, +0.029]) & 359 & 0.465 & 0.551 & 0.88 & $-$0.002 [$-$0.069, +0.061] \\
\bottomrule
\end{tabular}
}
\end{table}

\paragraph{Ranking or flag rate.} A coupling factor can raise the MCC at the
rule's constant in two ways: by ranking the god units higher, or by flagging
more units along much the same ranking. \Xs$\cdot$\Oflow{} flags
\matchedOOwnFlags{} of the reviewed confirm classes where \Xs{} flags
\matchedXOwnFlags{}. Two reads separate the two, post hoc
(Tables~\ref{tab:matched-class} and \ref{tab:matched-method}): the
threshold-free pooled AUC of each form's per-repository modified z over the
reviewed units, with the paired bootstrap against \Xs{}; and each form at a
matched flag rate, the constant that gives it as many flags over every unit
of the population's repositories as \Xs$\cdot$\Oflow{} has at 3.5. At class
level the outgoing factor ranks Blobs no better than complexity alone: AUC
\matchedOAucConfirm{} against \matchedXAucConfirm{} on confirm
(\matchedOAucVsXConfirm{}) and \matchedOAucExplore{} against
\matchedXAucExplore{} on explore (\matchedOAucVsXExplore{}); and at the
matched rate \Xs{} alone reaches \matchedXMatchedMcc{} (recall
\matchedXMatchedRecall{}) against the outgoing factor's \matchedOMcc{},
\matchedOVsXConfirm{} on confirm and \matchedOVsXExplore{} on explore. The
size null's AUC equals \Xs{}'s (\matchedNAucVsXConfirm{}), the divergence
factor's is resolved worse (\matchedDAucVsXConfirm{}), and the exploratory
size $\times$ calls form ranks no better either
(\matchedSizeCallsAucVsXConfirm{} on confirm,
\matchedSizeCallsAucVsXExplore{} on explore). On Fontana the outgoing factor
does rank better (\matchedFontanaOAucVsX{}; at the matched rate
\matchedFontanaOVsX{}). So on MLCQ the construct's one surviving reading at
the operating point, that outgoing data refines complexity, is a flag-rate
effect, and a lower constant on complexity alone does the same. At method level the ranking belongs to length. The any-call gate ranks
Long Methods better than the registered signal, cognitive complexity (AUC
\matchedGateAucConfirm{} against \matchedCogAucConfirm{},
\matchedGateAucVsXmConfirm{} on confirm and \matchedGateAucVsXmExplore{}
on explore), but no better than length itself: against tokens its
percentile-rank AUC differs by \matchedGatePctVsTokensConfirm{} on confirm
and \matchedGatePctVsTokensExplore{} on explore, and at tokens' flag rate
its MCC is \matchedGateMatchedVsTokensConfirm{} on confirm and
\matchedGateMatchedVsTokensExplore{} on explore. Lines and tokens rank
alike (percentile AUC \matchedLinesPctConfirm{} and
\matchedTokensPctConfirm{}, against \matchedCogPctConfirm{} for cognitive
complexity); the internal-call gate ranks worse
(\matchedIntGateAucVsXmConfirm{}), because the leaves it discards include
Long Methods; and the weight does not change the ranking at all
(\matchedWeightAucVsXmConfirm{}), its loss being pure flag rate. So the
registered signal was the wrong length measure, the registered factor was a
flag-rate loss on it, and the gate is a flag-rate gain on the right one.

\begin{table}[t]
\centering
\caption{Exploratory, post hoc: MLCQ's individual Long Method reviews on
measured units (\revLongReviews{} reviews by \revReviewers{} reviewers,
both halves), by the method's length and whether it is a leaf (no internal
call): the share of reviews that call the method a Long Method at any
severity.}
\label{tab:reviewers}
\small
\setlength{\tabcolsep}{4pt}
\fitwidth{% generated by paper/scripts/tables.py; do not edit
\begin{tabular}{lrrrr}
\toprule
Lines of the method & Leaf reviews & called Long Method & Non-leaf reviews & called Long Method \\
\midrule
1-10 & 770 & 0.4\% & 823 & 1.0\% \\
11-25 & 230 & 17.4\% & 571 & 24.3\% \\
26-50 & 80 & 80.0\% & 365 & 81.6\% \\
51-100 & 17 & 94.1\% & 151 & 94.0\% \\
101+ & 7 & 85.7\% & 40 & 95.0\% \\
\bottomrule
\end{tabular}
}
\end{table}

\paragraph{What the reviewers judge.} MLCQ keeps every review, so the
reading that a reviewer asks first whether a method is a self-contained
algorithm and then whether it is long can be put to the reviews themselves
(Table~\ref{tab:reviewers}; \revLongUnmapped{} reviews on units the harness
did not measure are left out). The share of reviews calling a method a
Long Method is \revNonLeafTenRate{} up to 10 lines,
\revNonLeafTwentyFiveRate{} at 11 to 25, \revNonLeafFiftyRate{} at 26 to 50
and \revNonLeafHundredRate{} at 51 to 100: a step at about twenty-five
lines. Each of the \revCutoffReviewers{} reviewers with ten or more
positives has a best length cutoff between \revCutoffMin{} and
\revCutoffMax{} lines (median \revCutoffMedian{}), separating their own
positive from their negative reviews at an MCC of \revCutoffMccMin{} to
\revCutoffMccMax{}; the median length of a reviewer's positives is
\revPosMedianMin{} to \revPosMedianMax{} lines and of their negatives at
most \revNegMedianMax{}. That is an absolute constant, carried from
reviewer to reviewer within a few lines, which is why a cutoff tuned on
MLCQ beat the per-repository rule in the companion study and why no
per-repository form here can reach it. At equal length a leaf is called a
Long Method no less often than a non-leaf (\revLeafFiftyRate{} against
\revNonLeafFiftyRate{} at 26 to 50 lines, \revLeafHundredRate{} against
\revNonLeafHundredRate{} at 51 to 100; over all long methods
\revLongLeafMinusNonLeaf{}), reviewers disagree on long leaves no more than
on long non-leaves (\revDisagreeLongLeaf{} against
\revDisagreeLongNonLeaf{} of multiply reviewed samples not unanimous), and
the severities match. The exemption the registered diagnostic supported is
narrower than the construct: among the reviewed methods cognitive
complexity flags, leaves are called Long Method less often than the rest
(\revFlaggedLeafRate{} against \revFlaggedRestRate{}) at the same median
length (\revFlaggedLeafMedianLen{} against \revFlaggedRestMedianLen{}
lines), so it is the branchy leaf, the dense algorithm complexity singles
out, that reviewers decline to call a Long Method, and not a leaf as such.
For Blob the reviews say less: the Blob rate is \revBlobHundredRate{} for
classes up to 100 lines, \revBlobTwoFiftyRate{} at 101 to 250,
\revBlobFiveHundredRate{} at 251 to 500 and then no higher
(\revBlobThousandRate{} at 501 to 1000, \revBlobOverRate{} above);
between \revBlobHundredDisagree{} and \revBlobTwoFiftyDisagree{} of
multiply reviewed samples are not unanimous; and among classes over 500
lines the Blob rate by tertile of leaving calls is
\revBlobLargeCallsLowRate{}, \revBlobLargeCallsMidRate{} and
\revBlobLargeCallsHighRate{}. The Blob label has a size step and much
disagreement, and little else these measures see, which is the ceiling
every class form hit.

\paragraph{A constant on length, and what was tried against it.} The
constant the reviews imply was read once as the plan reads a cutoff: chosen
on the explore half among absolute cutoffs on lines and read on the confirm
half (\lwConfirmRepos{} repositories, \lwConfirmUnits{} reviewed methods,
\lwConfirmPositives{} positives). The cutoff chosen is \lwLengthCutoff{}
lines or more. On confirm it flags \lwLengthFlagged{} methods with
\lwLengthHits{} hits, MCC \lwLengthMccConfirm{} (precision
\lwLengthPrecision{}, recall \lwLengthRecall{}), \lwLengthVsRule{} over
the best per-repository form, the registered rule on length gated on any
internal call; the MCC is flat from 20 to 30 lines (\lwNlocTwentyMcc{} at
20, \lwNlocThirtyMcc{} at 30) and has fallen by 40 lines
(\lwNlocFortyMcc{}), so the constant is a mid-twenties one rather than a
sharp number. Two further post hoc readings were then made on the same
half, each labelled as such, to ask whether anything a method's body
contains refines that constant. The first scored \lwHeldOutForms{} forms
built from the method's variables, the live width a reader holds at a line
in several definitions, the counts of its locals, parameters and fields,
and the flow of data from what enters the method to the calls it makes, by
the same procedure; every one is resolved below length
(\lwFormsBelowLength{} of \lwHeldOutForms{}), the best of them at
\lwBestFormMcc{}, \lwBestFormVsLength{}. The second scored the cost of
cutting a long method into pieces under the constant, the parameters and
return values the pieces would exchange; its cheapest-seam form reaches
\splitCostMcc{}, \splitCostVsLength{} against the constant, and among the
\splitWithinN{} long methods, the only ground such a form can gain, it
ranks the \splitWithinPositives{} positives worse than length does (AUC
\splitCostAucWithin{} against \splitNlocAucWithin{}, \splitCostAucVsNloc{}).
The share of a class's pieces bound to the class, read as a god-class
form, is last on every class oracle (\splitLockedShareMcc{} against
\splitBlobLengthMcc{} for the longest method on the confirm half;
\splitFontanaLockedMcc{} against \splitFontanaHvolMcc{} for Halstead
volume on Fontana, \splitFontanaLockedVsHvol{}). Nothing in a long method's
body, as these measures see it, says which of the methods length flags the
reviewers passed.

\paragraph{The logarithmic scale.} Every registered form was also recomputed
with $\ln(1 + X)$ in place of the complexity signal X (\logClassForms{}
class forms, \logMethodForms{} method forms), to ask whether the complexity
term dominates the product. Under the fixed constant a logarithm compresses
the tail and the rule then flags almost nothing on the signal alone
(\logXFlagsLog{} classes against \logXFlagsRaw{} for \Xs{}, MCC
\logXMccLog{}), so a log form recovers only where a factor re-expands the
tail; across the matrix the log form is resolved worse than the raw one in
\logClassWorse{} of \logClassForms{} class forms and better in
\logClassBetter{}, and the best log form against Blob reaches
\logClassTopMcc{}, no better than the best raw one. A monotone transform
moves the constant, not the ranking, which is why the threshold-free read
above is the one that decides.

\paragraph{What this is and is not.} These are the best of
\expClassForms{} class forms and \expMethodForms{} method forms, chosen by
looking, on labels that had all been read. The class-level gain over complexity alone is about +0.03 on each half,
unresolved, and absent from the threshold-free ranking; the method-level
gain is a recalibration of the constant on length, which ranks no better
than length itself. What the reviews show, that Long Method is an absolute
length, is the one reading here that does not depend on a chosen form, and
the constant read once on the confirm half is the one number here that was
chosen on explore and read as the plan reads a cutoff. Both are hypotheses for the
next registration, stated in Section~\ref{sec:disc}, and neither is a result
of this study.

\section{Deviations from the plan}
\label{sec:deviations}

The full account of every departure from the plan, with the numbers of
every analysis set aside, is in the ancillary data; this section states
each in brief.

\paragraph{The human populations were first scored on the repositories
with a positive only.} The plan scores each human oracle on every reviewed
unit, a repository whose reviewed units are all negatives included. The
harness attached a human oracle to a repository only when the repository
held a positive-labelled unit, so such repositories never entered the
population. The fault was found after a confirmatory analysis had been
written, its digest committed, and its results read and reported to the
author in a draft of this paper; that is the one departure from the plan's
order, and it is recorded in the plan as Amendment 10 with what had been
seen. The fix has no free parameter: it attaches the oracle whenever the
reviewed universe is present, which is the population the plan's words
define, and the units it restores are negatives that either arm may flag.
The explore-half choices and the confirmatory analysis were then made
again by the registered procedure, and the earlier files were set aside,
not deleted, with their digests. The fix restored \restoredBlobRepos{}
confirm repositories and \restoredBlobUnits{} reviewed classes to the Blob
population, \restoredLongRepos{} repositories and \restoredLongUnits{}
reviewed methods to Long Method, \restoredExploreBlobRepos{} repositories
to the explore half on which \Xs{} and the gate are chosen, and
\restoredFontanaSystems{} systems and \restoredFontanaUnits{} labelled
classes to Fontana's God Class. \Xs{} is \xClassChanged{} by it
(\sig{\priorXClass} before, \sig{\xClass} after) and \Xm{} is
\xMethodChanged{}; the gate moved from \priorGateRequiredExplore{} to
\gateRequired{}. On the set-aside analysis the primary read
\priorPrimMicro{}, \priorPrimVerdict{}, and the construct's verdict was
\priorConsVerdict{}; the ancillary data tabulates every headline reading
before and after.

\paragraph{The confirmatory analysis was written three times.} The first,
set aside by Amendment 9 before any of its results was read, had Fontana's
systems resolved for \fontanaVOneSystems{} of \fontanaSystems{} and read
the Fontana paired difference on \fontanaVOneUniv{} classes as
\fontanaVOneMicro{}, \fontanaVOneReading{}. The second is the one the
population fault set aside. The third, on the complete samples and the
registered populations, is the analysis this paper reports. All three are
in the ancillary data with their digests as committed.

\paragraph{Amendments before any confirm measurement.} Nine amendments
followed registration, each committed with a statement of what had been
seen, none after a confirm or seed-3 repository had been measured. Those
that change what is reported here: constructors are not inputs of a class
(the registered variant is kept beside the primary); the exemption-only and
amplification-only arms, the size null and the reading of the mechanism
were added on a review of the plan and the harness; the gate is sized from
the estimand's own bootstrap half-width, with the form first registered
reported beside it; the rule fitted on the non-exempt units, the
coupled-only read-out and the share of outgoing coupling from thrown
exceptions were added as sensitivities and diagnostics; the file rule is
matched against the source root; the measurer's budget and heap were set.
Faults in measuring the registered samples, each in the harness or the
machine and not in the design, were fixed and the affected repositories
measured again; the ancillary data lists them.

\paragraph{The gate was not met.} The confirm half carries
\primGateReposWithBlob{} repositories with a reviewed Blob against the
\gateRequired{} the explore half's half-width implied, and MLCQ has no more.
The primary is marked underpowered by the plan's rule, and its confirm
half-width is \primHalfWidthConfirm{}. What the gate says is that a
difference of 0.10 could have been missed on the primary, not that the
secondaries resolved on the same population are suspect.

\section{Discussion}
\label{sec:disc}

\paragraph{The construct as registered is not supported, and the data say
where it fails.} Three things were expected of the divergence factor: that
it would raise precision by setting complex leaves aside, that the leaves
set aside would be called god units less often, and that what a unit
receives would matter as well as what it sends. The second is seen, in two human oracles of three: a flagged catchment is
a Long Method, or a Fontana God Class, less often than a flagged junction,
and for MLCQ's Blob the difference is unresolved with every point estimate
on the same side; the individual reviews narrow it to the branchy leaves
complexity flags, and show no exemption for leaves at equal length. The
first and the third are not seen. The precision falls and the recall rises,
and in every family of every class-level oracle, subtracting the inputs from
the outputs costs agreement. What is seen at the registered operating point is that the outputs alone
carry a signal, resolved against both human class oracles and every
class-level tool oracle; the threshold-free read (Section~\ref{sec:explore-forms})
then shows that on MLCQ the outgoing factor ranks Blobs no better than
complexity alone and gains by flagging two and a half times as many
classes, while on Fontana it does rank better. The verdict \emph{outgoing
data only} is therefore a statement about the operating point, and the
construct's residue on the human labels is smaller than the verdict reads.
The size null says how far the operating-point signal is size: a plain weighting by the number of methods beats the
divergence factor by a resolved margin and is itself not resolved against
complexity alone. The outgoing factor's effect on \Xs{} is of the same
order as the size null's at the operating point and larger in the family
averages, so part of it is size and part may not be; separating the two is
the next question, and this study did not register a test of it. In Tempero
and Ralph's terms \cite{tempero2026metrics}, divergence as an indicator of
the sinew construct fails convergent validity against three human
definitions of a god unit, and the component that converges, at the operating point, is out-coupling,
which Arif, Kuutila and Ralph found loading as a construct of its own
\cite{arif2026}, and which the companion study's two readers agreed was what
distinguished the extra flags they could read \cite{paperone}.

\paragraph{The amplification, not the exemption, is where the factor
fails.} At class level the whole
factor is resolved worse than the exemption alone and not better than the
amplification alone, and at method level the amplification alone accounts
for the whole loss while the exemption alone is within hundredths of
complexity. A factor of $d + 1$ is
unbounded and, for a class, grows with the number of calls the class makes,
so the modified z of X$\cdot$D is dominated by the classes that call most,
whatever their complexity. The outgoing factor has the same unbounded
growth and does better, which says the sign of the inputs is the problem:
a class with a wide callable surface is one that other code uses, and
subtracting its surface pushes the hubs people call Blobs towards the
catchment side. That is the opposite of the construct's reading of a
catchment as a leaf doing honest work.

\paragraph{Where the exemption seemed to help, and what it was.} The leaf
diagnostic says that among the units complexity flags, a catchment is a god
unit less often by two human oracles, and the reviews confirm it for the
dense, branchy leaves cognitive complexity singles out. But the
exemption-only arm is within hundredths of complexity alone at both levels,
because the leaves it sets aside include god units and the zeros it leaves
move the rule's scale, and the whole factor then amplifies the junctions
that remain by an unbounded $d + 1$ and loses what the exemption gained.
The post hoc gate on any call, which looked like the exemption's own form,
turned out to be something else: methods that call nothing are almost never
long, so the gate exempts no Long Method; it zeros the short trivial
methods and so moves the constant to where length separates best. It ranks
Long Methods no better than tokens or lines do. The finding underneath is
that length ranks Long Methods best of everything measured, better than the
registered cognitive signal, and that the rule's constant on raw length is
badly placed: a tuned constant on length reached 0.785 in the companion
study and \lwLengthMccConfirm{} at \lwLengthCutoff{} lines here, and the
reviews explain why, since each reviewer carries a length cutoff near
twenty-five lines whatever the repository.

\paragraph{Two hypotheses for the next registration.} The exploration leaves
two forms, each old in its ingredients. A god class is a class whose total
lines times its outgoing calls is an outlier for its codebase: size times
fan-out, which is Chidamber and Kemerer's response for a class as a product
rather than a sum, and Lanza and Marinescu's God Class without the cohesion
clause. It is ahead of complexity alone by about +0.03 on each half, unresolved,
ahead of a size-only weighting on one half, and no better than complexity
alone in the threshold-free ranking; the next study's
primary is that form against its size-only counterpart, and against CBO and
RFC, on labels no one has read. At method level the hypothesis is not a coupling form at all. What
MLCQ's reviewers label a Long Method is a method of \lwLengthCutoff{}
lines or more, leaf or not; the right detector for that label is a
constant on length, which the companion study's tuned cutoff already
reaches, and the per-repository rule is the wrong model for it. Whatever
else the term god method means, the registered construct did not capture
it, and nor did the live width, the variable counts, the data flow or the
split cost read after it. What remains to register at method level is
whether the branchy-leaf exemption the reviews do show, a dense leaf
complexity flags that reviewers decline to call a Long Method, adds
anything to a constant on length, and whether the per-repository rule has
a place at method level on code whose scale has moved, which is the purpose
the plan stated and this data cannot test. The construct's next form at
method level is not a detector but a cost: with the trigger fixed at the
constant, the data that crosses the seam when a long method is cut into
pieces under it, the sinew of the refactoring rather than of the method.
Prior tools rank extract-method candidates within a method
\cite{tsantalis2011, haas2016}; none takes the length cap as a constraint
with a cost on what the pieces exchange, and that instrument is to be built
and read descriptively before any claim is registered for it.
Both hypotheses are post hoc and were found by searching; the power
estimate from this study's intervals says the class-level claim needs
either several times the labelled repositories or a better measure.

\paragraph{The human Blob is weakly predicted by any of this.} Complexity
alone reaches an MCC of \primXMcc{} against the reviewers on the confirm
half, the best arm in the paper, the outgoing factor, \secOFlowMcc{},
and the companion
study's exploratory comparison found no tool definition above about 0.2 on
the full population \cite{paperone}; classifiers trained on the labels reach
0.51 \cite{mlcq2023}. The reviewers labelled without a given definition and
agree with each other moderately \cite{schumacher2010}, so part of the
ceiling is the oracle. The differences read here are differences between
arms on the same units and the same repositories, and the paired design
removes the oracle's noise from the comparison but not from the ceiling.

\paragraph{Calls the parser cannot see.} A quarter of calls on the confirm
half resolve to nothing, and injected collaborators are invisible, so
$\mathit{out}$ reads low and some hubs look like leaves. Removing the
repositories that declare a DI framework moved the primary against the
factor, which does not support reading the result as an artefact of the
missing calls, but the sensitivity removes repositories, not the missing
calls within the rest.

\section{Threats to validity}
\label{sec:threats}

\paragraph{Construct.} Divergence counts parameters and arguments, not
what travels in them; a request object is one parameter carrying many
items, and a class whose data arrives bundled reads as receiving little.
Return values and exceptions are not counted. A class's inputs count the
parameters of its callable methods, so a wide interface with many small
methods reads as receiving much, which the discussion identifies as the
likely failure of the sign; a different definition of a class's inputs is a
different construct and was not registered.

\paragraph{Internal.} The populations are the registered ones only because
a fault in the harness was fixed and the analysis made again after its
first results had been read (Section~\ref{sec:deviations}); the fix had no
free parameter, and the ancillary data gives every headline reading before
and after it. The two arms
of the primary are fitted under different scale estimators because the
weighted form is zero on most units; the sensitivity that removes the
difference brings the arm to complexity alone. Unresolved calls make
$\mathit{out}$ a lower bound. The complexity extractor and the measurer
disagree on cognitive complexity for about one method in twenty; the
JavaParser-only set carries no such disagreement and agrees on every
verdict.

\paragraph{External.} Java only; a dynamically typed language resolves no
call statically. MLCQ's repositories are concentrated in a few organisations
(\mlcqOwnerTopShare{} of the confirm half in ten owners,
\mlcqApacheRepos{} Apache) and the seed-3 sample is not
(\seedThreeOwnerTopShare{}). Fontana's instances are extreme. The purpose
stated in the plan, a definition that moves with the codebase as
agent-maintained code changes what a normal level of complexity is, was not
tested: no agent-written corpus with human labels exists at the scale a
pre-registered test needs, and the companion study's second reading
applies, that a detector must first match the single scale of existing
human-written code.

\paragraph{Statistical.} The primary is underpowered by the plan's own
gate, and no secondary is corrected for multiplicity; the registered
readings are stated as such and the rest labelled exploratory. The
consistency analysis's cells share one bootstrap, so the count of resolved
cells is a description, not a count of independent tests. The intervals are
percentile bootstraps over repositories, which with \primRepos{}
repositories and a few reviewed units each are coarse at the tails. The
post hoc readings of Section~\ref{sec:explore-forms} were made on the
confirm half after the constant on length had been chosen on explore and
read on it, so that half is a hold-out for the constant alone and not for
anything chosen after it.

\section{Reproducibility and data availability}
\label{sec:repro}

Every number in this paper is produced by a script from the harness's
confirmatory analysis file, and every table is a generated fragment. The
ancillary data holds the confirmatory analysis and the two that were set
aside, with their digests as committed, and the full account of every
deviation from the plan; the explore-half decisions and the
measurer's validation; the per-set run records and diagnostics; the sample
files that pin every repository and system to its commit or archive; the
converted MLCQ labels and the Fontana labels with the record of where the
corpus came from and the maintainer's reply; and the list of MLCQ
repositories whose labelled commits moved. The post hoc exploration's scripts, their outputs and their digests are
there too, labelled as such. The raw per-repository records regenerate by running the
harness on the sample files. The ancillary README
names the public repositories of the harness, the companion study, the
sampler and the measurer, the commit of each, and the command behind every
number.

\section{Conclusion}

Weighting a unit's complexity by its coupling divergence does not recover
the classes and methods people call god units better than complexity
alone. Against MLCQ's reviewers the primary resolves no difference and is
underpowered, and every resolved reading of the whole factor goes the other
way: it loses to the exemption alone, to a weighting by the number of
methods, to complexity alone on the threshold-free ranking, and, without
calls to helpers, to complexity alone at method level outright. The
exemption the claim rests on is partly in the labels: a complex method that
receives more than it sends is called a Long Method less often, and such a
class a God Class less often by Fontana's team, while MLCQ's reviewers call
it a Blob not detectably less often. What the factor adds to the exemption,
an amplification by the divergence of a junction, costs more than the
exemption gains. At the registered operating point the one component that carries a signal
is the data a class sends out, resolved against both human class oracles
and every class-level tool oracle, and subtracting what the class receives
costs agreement in every family; the construct's verdict, by its
pre-registered rule, is outgoing data only. Read threshold-free, that gain
is a flag rate on MLCQ and a better ranking only on Fontana, so a god class
is a class that sends much out only in the sense that flagging more large
classes finds more Blobs.
How much of that is the size of the class, and what remains of coupling
once size is held, are the questions this result leaves for a registered
test, with the post hoc exploration's reading as the hypothesis at class level,
total lines times outgoing calls against its size-only counterpart. At
method level the reviews settle more than the forms do: MLCQ's reviewers
call a method a Long Method past about twenty-five lines, each within a few
lines of the others, leaf or not; length ranks their verdicts best of
everything measured; and a constant on length, \lwLengthCutoff{} lines or
more at MCC \lwLengthMccConfirm{} on the confirm half, not a
per-repository outlier rule, is the detector that matches them. Nothing
read after it in a long method's body, its live width, its variables, its
data flow or the cost of cutting it, refines that constant; the cost of the
cut is the form of sinew left to build. The per-repository rule's
case rests on code whose scale has moved, and that is the test still to
run.

\section*{Competing interests}

The author created and maintains OfficeFloor, a framework whose design
rests on taking data off the control flow, and the sinew construct that
divergence was meant to indicate comes from that line of work; a positive
result would have supported the framework's premise. The result is
negative on the registered claim, and the reading that survives, that
outgoing data carries the signal, was recorded as a competing expectation
before the data existed, from the author's own exploratory runs. The
fault in the harness's populations (Section~\ref{sec:deviations}) was
found after the confirmatory results had been read; the fix is the one the
plan's text dictates, and every headline reading before and after it is
reported; the plan, the seed, the decision rule
and the gate were fixed before the data existed, and a replication needs no
cooperation from the author. The companion study, the sampler and the
measurer are the author's.

\section*{Funding}

This work received no external funding and was carried out and paid for by
the author.

\begin{thebibliography}{99}

\bibitem{arif2026} M.~Arif, M.~Kuutila and P.~Ralph,
  ``A factor analysis of class-level software metrics,''
  in \emph{Proc. ACM/IEEE International Symposium on Empirical Software
  Engineering and Measurement (ESEM)}, 2026; arXiv:2609.20411.

\bibitem{fontana2016} F.~Arcelli Fontana, M.~V. M\"antyl\"a, M.~Zanoni and
  A.~Marino,
  ``Comparing and experimenting machine learning techniques for code smell
  detection,''
  \emph{Empirical Software Engineering}, vol.~21, no.~3, pp.~1143--1191, 2016.

\bibitem{cktool} M.~Aniche,
  ``Java code metrics calculator (CK),'' 2015. Version 0.7.0.

\bibitem{ck1994} S.~R. Chidamber and C.~F. Kemerer,
  ``A Metrics Suite for Object Oriented Design,''
  \emph{IEEE Transactions on Software Engineering}, vol.~20, no.~6,
  pp.~476--493, 1994.

\bibitem{campbell2018} G.~A. Campbell,
  ``Cognitive Complexity: a new way of measuring understandability,''
  SonarSource white paper, 2018.

\bibitem{dinucci2018} D.~Di~Nucci, F.~Palomba, D.~A. Tamburri, A.~Serebrenik
  and A.~De~Lucia,
  ``Detecting code smells using machine learning techniques: Are we there
  yet?''
  in \emph{Proc. 25th IEEE International Conference on Software Analysis,
  Evolution and Reengineering (SANER)}, pp.~612--621, 2018.

\bibitem{elemam2001} K.~El~Emam, S.~Benlarbi, N.~Goel and S.~N. Rai,
  ``The confounding effect of class size on the validity of object-oriented
  metrics,''
  \emph{IEEE Transactions on Software Engineering}, vol.~27, no.~7,
  pp.~630--650, 2001.

\bibitem{friedman1937} M.~Friedman,
  ``The use of ranks to avoid the assumption of normality implicit in the
  analysis of variance,''
  \emph{Journal of the American Statistical Association}, vol.~32, no.~200,
  pp.~675--701, 1937.

\bibitem{haas2016} R.~Haas and B.~Hummel,
  ``Deriving extract method refactoring suggestions for long methods,'' in
  \emph{Software Quality Days (SWQD 2016)}, LNBIP vol.~238, Springer, 2016,
  pp.~144--155.

\bibitem{henry1981} S.~Henry and D.~Kafura,
  ``Software structure metrics based on information flow,''
  \emph{IEEE Transactions on Software Engineering}, vol.~SE-7, no.~5,
  pp.~510--518, 1981.

\bibitem{iglewicz1993} B.~Iglewicz and D.~C. Hoaglin,
  \emph{How to Detect and Handle Outliers}. ASQC Quality Press, 1993.

\bibitem{kitchenham1990} B.~A. Kitchenham, L.~M. Pickard and S.~J. Linkman,
  ``An evaluation of some design metrics,''
  \emph{Software Engineering Journal}, vol.~5, no.~1, pp.~50--58, 1990.

\bibitem{lanza2006} M.~Lanza and R.~Marinescu,
  \emph{Object-Oriented Metrics in Practice}. Springer, 2006.

\bibitem{mlcq2020} L.~Madeyski and T.~Lewowski,
  ``MLCQ: Industry-relevant code smell data set,''
  in \emph{Proc. 24th International Conference on Evaluation and Assessment in
  Software Engineering (EASE)}, pp.~342--347, 2020. Data set version 1.1,
  Zenodo record 3666840.

\bibitem{mlcq2023} L.~Madeyski and T.~Lewowski,
  ``Detecting code smells using industry-relevant data,''
  \emph{Information and Software Technology}, vol.~155, 107112, 2023.

\bibitem{marinescu2004} R.~Marinescu,
  ``Detection strategies: metrics-based rules for detecting design flaws,''
  in \emph{Proc. 20th IEEE International Conference on Software Maintenance
  (ICSM)}, pp.~350--359, 2004.

\bibitem{mcCabe1976} T.~J. McCabe,
  ``A Complexity Measure,''
  \emph{IEEE Transactions on Software Engineering}, vol.~SE-2, no.~4,
  pp.~308--320, 1976.

\bibitem{moha2010} N.~Moha, Y.-G. Gu\'eh\'eneuc, L.~Duchien and A.-F. Le~Meur,
  ``DECOR: A Method for the Specification and Detection of Code and Design
  Smells,''
  \emph{IEEE Transactions on Software Engineering}, vol.~36, no.~1, pp.~20--36,
  2010.

\bibitem{paperone} D.~Sagenschneider,
  ``Complexity concentration as a god-class and god-method detector: a
  pre-registered test of a scale-free outlier rule against tool definitions
  and human labels,'' 2026. \url{https://doi.org/10.5281/zenodo.23581333}

\bibitem{schumacher2010} J.~Schumacher, N.~Zazworka, F.~Shull, C.~Seaman and
  M.~Shaw,
  ``Building empirical support for automated code smell detection,''
  in \emph{Proc. ACM-IEEE International Symposium on Empirical Software
  Engineering and Measurement (ESEM)}, article 8, 2010.

\bibitem{tempero2010} E.~Tempero, C.~Anslow, J.~Dietrich, T.~Han, J.~Li,
  M.~Lumpe, H.~Melton and J.~Noble,
  ``The Qualitas Corpus: A curated collection of Java code for empirical
  studies,''
  in \emph{Proc. 17th Asia Pacific Software Engineering Conference (APSEC)},
  pp.~336--345, 2010.

\bibitem{tempero2026metrics} E.~Tempero and P.~Ralph,
  ``Making software metrics useful,''
  \emph{IEEE Computer}, in press, 2026; arXiv:2603.16012.

\bibitem{tsantalis2011} N.~Tsantalis and A.~Chatzigeorgiou,
  ``Identification of extract method refactoring opportunities for the
  decomposition of methods,'' \emph{Journal of Systems and Software},
  vol.~84, no.~10, pp.~1757--1782, 2011.

\bibitem{yourdon1979} E.~Yourdon and L.~L. Constantine,
  \emph{Structured Design: Fundamentals of a Discipline of Computer Program
  and Systems Design}. Prentice-Hall, 1979.

\bibitem{zhang2025deps} Z.~Zhang, E.~Wen and E.~Tempero,
  ``Analyzing variations in dependency distributions due to code smell
  interactions,'' arXiv:2509.03896, 2025.

\end{thebibliography}

\end{document}
